\documentclass[11pt]{article}

\usepackage[utf8]{inputenc}
\usepackage[T1]{fontenc}
\usepackage{natbib}
\usepackage{geometry}
\geometry{margin=2.5cm}
\usepackage{hyperref}
\usepackage{url}
\usepackage{booktabs}
\usepackage{amsmath}
\usepackage{amsfonts}
\usepackage{amssymb}
\usepackage{nicefrac}
\usepackage{microtype}
\usepackage{xcolor}
\usepackage{graphicx}
\usepackage{float}
\usepackage[table]{xcolor}
\usepackage{subcaption}
\usepackage{cleveref}
\crefname{equation}{Eq.\@}{Eqs.\@}
\Crefname{equation}{Eq.\@}{Eqs.\@}
\crefname{figure}{Fig.\@}{Figs.\@}
\Crefname{figure}{Fig.\@}{Figs.\@}
\crefname{table}{Tab.\@}{Tabs.\@}
\Crefname{table}{Tab.\@}{Tabs.\@}
\crefname{section}{Sect.\@}{Sects.\@}
\Crefname{section}{Sect.\@}{Sects.\@}

\definecolor{cgreen}{HTML}{2ea043}
\definecolor{corange}{HTML}{d29922}
\definecolor{cred}{HTML}{cf222e}
\newcommand{\good}[1]{\textcolor{green!50!black}{#1}}
\newcommand{\bad}[1]{\textcolor{orange}{#1}}

\def\eg{{\it e.g.}}
\def\ie{{\it i.e.}}
\def\R{\mathbb{R}}
\def\vrest{V^{\mathrm{rest}}}
\def\relu{\mathrm{ReLU}}
\def\softplus{\mathrm{softplus}}
\def\hulse{Hulse et al.\ 2021}

% Functional-class macros mirroring zebrafish.tex (R/L/D/Z analysis)
\newcommand{\mi}{\mathcal{I}}
\newcommand{\omag}{\mathcal{W}}
\newcommand{\omod}{\mathcal{O}}
\newcommand{\specf}{\mathcal{T}}
\newcommand{\cls}[1]{\mathsf{#1}}

\title{\textbf{\emph{Drosophila} central complex: identifiable computation, non-identifiable implementation}}
\author{}
\date{}

\begin{document}
\maketitle
\begin{abstract}
Theoretical work on emergent structure in recurrent neural networks
almost never has a ground-truth comparison: discovered modules are
evaluated against themselves. The \emph{Drosophila} central complex
(CX) is one of the rare exceptions --- a complete connectome
\citep{hulse2021} pairs with a known biological function
(head-direction integration) and a published RNN model
\citep{vafidis2022}. We use this pairing to ask what optimisation of the path-integration
task identifies when the connectome topology is given. Four constrained
parameterisations are trained on the same input--output map under
10-fold cross-validation, with and without a total-variation
regulariser on the EPG ring:
a Known-ODE RNN, a frozen-connectome control, a fully connected RNN,
and a message-passing GNN with learned neuronal signalling functions.
The frozen-connectome control fails to reach path-integration accuracy ($p_{\mathrm{acc}} = 0.29$, 0/10 folds above $|r| \ge 0.9$); the non-trivial residual $|r| = 0.669$ shows the raw connectome fixes topology while training supplies the integrator gain;
the fully connected RNN converges only intermittently (3/10); the Known-ODE RNN and
the GNN converge reliably (10/10 with TV). Across folds the learned
recurrent weights cluster more tightly with one another than with the
raw hemibrain template, indicating that training supplies edge
magnitudes the topology alone does not. The GNN reaches comparable
behavioural performance through a different implementation: its
learned drift and signalling functions deviate from the hardcoded
leak/sigmoid of the Known-ODE RNN. Thus optimising the
path-integration task recovers the ring-attractor computation, but
not a unique circuit implementation. Resolving these degeneracies will require a new methodology: looking for consistency between connectome, neural activity and behavioral task. 
\end{abstract}

\section{Introduction}


\section{Methods}
\label{sec:methods}

\textbf{Recurrent neural network.}\label{method:RNN}
The RNN model is a sign-locked discrete-time recurrent network whose connectivity is constrained to the \emph{Drosophila} hemibrain central-complex connectome. For each neuron $i\in\{1,\dots,N\}$ the subthreshold activity $\hat{h}_i(t)$ evolves according to
\begin{equation}
\tau\,\frac{d\hat{h}_i(t)}{dt} = -\hat{h}_i(t)
  + \sum_j \hat{W}^{\mathrm{rec}}_{ij}\,\sigma\!\big(\hat{h}_j(t)\big)
  + \sum_l \hat{W}^{\mathrm{in}}_{il}\,u_l(t)
  + \hat{b}_i,
\label{eq:cx_dynamics}
\end{equation}
\begin{equation}
\hat{y}(t) = \big(\cos\hat{\theta}(t),\,\sin\hat{\theta}(t)\big) = \hat{W}^{\mathrm{out}}\sigma\big(\hat{h}(t)\big) + \hat{b}^{\mathrm{out}},
\label{eq:cx_readout}
\end{equation}
where $\hat{y}(t)\in\mathbb{R}^2$ is the network's inference of the head-direction signal, encoded as the $(\cos\hat\theta,\sin\hat\theta)$ pair so that the supervisory loss matches the ground-truth heading on the unit circle and avoids the $\pm 2\pi$ wraparound; the decoded heading is recovered as $\hat\theta(t) = \mathrm{atan2}\big(\hat y_2(t), \hat y_1(t)\big)$. The activation $\sigma$ is the sigmoid, the membrane time constant is $\tau = 0.1\,\mathrm{s}$ and the integration step is $\Delta t = 0.01\,\mathrm{s}$. The input $u(t)\in\mathbb{R}^3 = (\omega(t),\,\cos\theta_0\,\delta_{t,0},\,\sin\theta_0\,\delta_{t,0})$ encodes angular velocity and a one-frame initial-heading impulse. The recurrent matrix is sign-locked to the connectome template $W^{\mathrm{con}}\in\mathbb{R}^{N\times N}$: the trainable per-edge parameter is the magnitude $\hat{S}\in\mathbb{R}^{N\times N}$, and
\begin{equation}
\hat{W}^{\mathrm{rec}}_{ij} = |\hat{S}_{ij}|\,\mathrm{sign}(W^{\mathrm{con}}_{ij})\,(1-\delta_{ij}).
\label{eq:cx_w_rec}
\end{equation}
The angular-velocity input is anatomically restricted to the PEN subpopulations through a strict four-scalar gate on the angular-velocity column of $\hat{W}^{\mathrm{in}}$,
\begin{equation}
\hat{W}^{\mathrm{in}}_{i,\omega} = v_s \quad \text{if } i \in s,\ \ s \in \{\mathrm{PENa}_L,\mathrm{PENa}_R,\mathrm{PENb}_L,\mathrm{PENb}_R\},
\label{eq:cx_pen_gate}
\end{equation}
with four learnable scalars $\{v_{\mathrm{PENa}_L},\,v_{\mathrm{PENa}_R},\,v_{\mathrm{PENb}_L},\,v_{\mathrm{PENb}_R}\}$ initialised antisymmetrically across hemispheres. The remaining two columns of $\hat{W}^{\mathrm{in}}$ --- carrying the initial-heading impulse $(\cos\theta_0,\sin\theta_0)\,\delta_{t,0}$ --- remain free for all neurons. The learnable parameters are
\begin{equation*}
\big\{\hat{S},\;
v_{\mathrm{PENa}_L},\,v_{\mathrm{PENa}_R},\,v_{\mathrm{PENb}_L},\,v_{\mathrm{PENb}_R},\;
\hat{W}^{\mathrm{in}}_{:,\theta_0},\,\hat{b},\,\hat{W}^{\mathrm{out}},\,\hat{b}^{\mathrm{out}}\big\};
\end{equation*}
the connectome template $W^{\mathrm{con}}$, PEN subpopulation membership, neuron-type assignments, $\tau$ and $\Delta t$ are given. Training optimises a per-frame MSE on the heading readout with two connectome-prior regularisers, an $\ell_1$ sparsity prior on the recurrent magnitudes, and a circular total-variation prior on the EPG ring firing rates,
\begin{equation}
\mathcal{L}
= \tfrac{1}{BT}\sum_{b,t} \|\hat{y}_{b,t} - y_{b,t}\|_2^2
+ \lambda_{\mathrm{cos}}\,\mathcal{L}_{\mathrm{cos}}
+ \lambda_{\mathrm{nf}}\,\mathcal{L}_{\mathrm{nf}}
+ \lambda_{1}\,\|\hat{S}\|_1
+ \lambda_{\mathrm{tv}}\,\mathcal{L}_{\mathrm{tv}},
\label{eq:cx_loss}
\end{equation}
with block losses computed over type-pair masks $\mathcal{B}_p\subset\{1,\dots,N\}^2$ (one block per ordered pre/post neuron-type pair $p\in\mathcal{P}$),
\begin{align}
\mathcal{L}_{\mathrm{cos}} &= \tfrac{1}{|\mathcal{P}|}\sum_{p\in\mathcal{P}}\!\left[1 - \frac{\langle \hat{W}^{\mathrm{rec}}_{\mathcal{B}_p},\,W^{\mathrm{con}}_{\mathcal{B}_p}\rangle}{\|\hat{W}^{\mathrm{rec}}_{\mathcal{B}_p}\|\,\|W^{\mathrm{con}}_{\mathcal{B}_p}\|}\right],
\label{eq:cx_cos}\\
\mathcal{L}_{\mathrm{nf}} &= \tfrac{1}{|\mathcal{P}|}\sum_{p\in\mathcal{P}} \relu\!\big(\kappa - \overline{|\hat{W}^{\mathrm{rec}}_{\mathcal{B}_p}|}\big)^2,
\label{eq:cx_nf}\\
\mathcal{L}_{\mathrm{tv}} &= \tfrac{1}{BT}\sum_{b,t}\,\sum_{i=1}^{|\mathcal{R}|}
\Big|\sigma\!\big(\hat h_{\pi_{\mathcal{R}}(i+1)}(t)\big)
   - \sigma\!\big(\hat h_{\pi_{\mathcal{R}}(i)}(t)\big)\Big|,
\label{eq:cx_tv}
\end{align}
where $\pi_{\mathcal{R}}$ is the EPG ring permutation that orders the EPG neurons by their PB glomerulus / preferred-phase index around the ellipsoid body, indices are interpreted modulo $|\mathcal{R}|$ to close the ring, and $\sigma$ is the sigmoid firing-rate non-linearity of \cref{eq:cx_dynamics}.
The cosine-distance term $\mathcal{L}_{\mathrm{cos}}$ anchors each block to the connectome template; the norm-floor term $\mathcal{L}_{\mathrm{nf}}$ prevents any type-pair block from collapsing to zero by penalising the (squared) shortfall of the block mean absolute weight $\overline{|\hat{W}^{\mathrm{rec}}_{\mathcal{B}_p}|}$ below a fixed floor $\kappa$. The threshold $\kappa$ is a hyperparameter (typical value $\kappa = 0.05$) that sets the minimum functional magnitude each type-pair block must retain; because $\mathcal{L}_{\mathrm{cos}}$ is scale-invariant, $\mathcal{L}_{\mathrm{nf}}$ is required to prevent a degenerate solution in which the cosine prior is satisfied by an arbitrarily small recurrent magnitude. The circular total-variation term $\mathcal{L}_{\mathrm{tv}}$ is a \emph{prior on the solution} rather than a constraint imposed by the path-integration task: the ring permutation $\pi_{\mathcal{R}}$ orders the 46 EPG nodes by their PB-glomerulus / EB-column index, so $\mathcal{L}_{\mathrm{tv}}$ penalises spatial jumps in the activity \emph{across EB columns} that the EPG arbors tile, not jumps across the cortical EPG somata (which sit outside the neuropil envelope and are not where the bump is measured). Calcium imaging records GCaMP fluorescence from EPG arbors in the EB and shows that the EPG population forms a \emph{single} coherent bump that moves smoothly around the EB as the fly turns \citep{seelig2015,kim2017,green2017}; $\mathcal{L}_{\mathrm{tv}}$ shapes the same EB-column profile that this imaging measures, propagating that anatomical-consistency prior to the rest of the network via the recurrent dynamics. A linear-reconstruction analysis of the null space of the trained $\hat{W}^{\mathrm{rec}}$ under the $+$TV training recipe is given in \cref{app:degeneracy_cx}.

\bigskip
\noindent\textbf{Graph neural network model.}\label{method:GNN}
The GNN model replaces the dense recurrent operator $\sum_j \hat{W}^{\mathrm{rec}}_{ij}\sigma(\hat{h}_j)$ of \cref{eq:cx_dynamics} with a message-passing graph neural network \citep{gilmer_neural_2017}, and --- crucially --- does \emph{not} hard-code either the leak, the firing-rate nonlinearity, or the synaptic activation function: the linear leak $-\hat{h}_i/\tau$, the sigmoid $\sigma$, and the per-edge message that combines a presynaptic activity with its postsynaptic effect must all be \emph{discovered} from the input--output map inside the learned drift $f_\theta$ and the learned signalling function $g_\phi$.
Beyond removing this functional bias, the GNN also admits an explicit form of neuronal heterogeneity through a per-neuron learnable embedding $\mathbf{a}_i\in\mathbb{R}^2$ that enters both $f_\theta$ and $g_\phi$: same-type neurons share their architectural role but can specialise via their embedding, so a single pair of MLPs covers the seven CX cell types and a hundred-plus instance-level distinctions without any further per-type parameter tying. Nodes correspond to neurons and edges to the non-zero entries of the hemibrain connectome template $W^{\mathrm{con}}$, and the subthreshold activity evolves as
\begin{equation}
\tau\frac{d\hat{h}_i(t)}{dt} =
  f_\theta\!\Big(\hat{h}_i(t),\,\mathbf{a}_i,\,
       \sum_{j\in\mathcal{N}_i} \hat{W}_{ij}\,g_\phi\!\big(\hat{h}_j(t),\mathbf{a}_j\big)\Big)
  + \sum_l \hat{W}^{\mathrm{in}}_{il}\,u_l(t),
\label{eq:cx_gnn_dynamics}
\end{equation}
with the same readout (\cref{eq:cx_readout}), PEN-gated encoder (\cref{eq:cx_pen_gate}), and dynamics constants ($\tau,\Delta t$) as the RNN. The neighbour set is fixed by the connectome topology, $\mathcal{N}_i = \{j : W^{\mathrm{con}}_{ij}\neq 0\}$. The node update $f_\theta$ and edge message $g_\phi$ are three-layer perceptrons (width 64, ReLU, linear output). $g_\phi$ maps the presynaptic state $\hat{h}_j$ to a per-edge message conditioned on the latent embedding $\mathbf{a}_j$ and weighted by the per-edge gain $\hat{W}_{ij}$; $f_\theta$ then processes the postsynaptic state $\hat{h}_i$, the embedding $\mathbf{a}_i$, and the aggregated synaptic drive. We constrain $g_\phi$ to be monotonically increasing in $\hat{h}_j$ via the prior $\partial g_\phi/\partial \hat{h} \ge 0$ (see \cref{eq:cx_gnn_loss}); combined with the sign-locked $\hat{W}_{ij}$ (\cref{eq:cx_gnn_w}) this restores Dale's law at the per-edge level without introducing a $\pm$-sign degeneracy. The leak, the firing-rate nonlinearity, and any departure from a linear synapse must emerge from $(f_\theta, g_\phi, \mathbf{a})$ themselves; the monotonicity priors below provide the minimum inductive bias for $f_\theta, g_\phi$ to behave like a neuron (postsynaptic drive increases with net excitatory input). The per-edge weight is sign-locked to the connectome as in the RNN,
\begin{equation}
\hat{W}_{ij} = |W_{ij}|\,\mathrm{sign}(W^{\mathrm{con}}_{ij})\,(1-\delta_{ij}),
\label{eq:cx_gnn_w}
\end{equation}
where $W_{ij}\in\mathbb{R}$ is the trainable per-edge magnitude. The learnable parameters are
\begin{equation*}
\big\{\theta,\,\phi,\,\mathbf{a},\,W,\;
v_{\mathrm{PENa}_L},\,v_{\mathrm{PENa}_R},\,v_{\mathrm{PENb}_L},\,v_{\mathrm{PENb}_R},\;
\hat{W}^{\mathrm{in}}_{:,\theta_0},\,\hat{b},\,\hat{W}^{\mathrm{out}},\,\hat{b}^{\mathrm{out}}\big\};
\end{equation*}
the connectome template $W^{\mathrm{con}}$, PEN subpopulation membership, neuron-type assignments, $\tau$, $\Delta t$, and ring orderings are given.

As with the RNN, heading recovery from cos/sin readouts is fundamentally non-unique: many distinct $(\theta,\phi,\mathbf{a},W)$ configurations produce identical readouts. We constrain the solution space with the same connectome-prior regularisers as the RNN ($\mathcal{L}_{\mathrm{cos}}$, $\mathcal{L}_{\mathrm{nf}}$ from \cref{eq:cx_cos,eq:cx_nf}, computed on the per-edge gain surface $\hat{W}_{ij}$ with floor threshold $\kappa$), the same circular-TV calcium-imaging prior $\mathcal{L}_{\mathrm{tv}}$ on the EPG ring firing rates (\cref{eq:cx_tv}), and sparsity and monotonicity priors on the MLPs:
\begin{align}
\mathcal{L} \;=\;& \tfrac{1}{BT}\sum_{b,t}\|\hat{y}_{b,t} - y_{b,t}\|_2^2
+ \lambda_{\mathrm{cos}}\mathcal{L}_{\mathrm{cos}}
+ \lambda_{\mathrm{nf}}\mathcal{L}_{\mathrm{nf}}
+ \lambda_{\mathrm{tv}}\mathcal{L}_{\mathrm{tv}} \nonumber\\
& +\;\lambda_0\|\theta\|_1 + \lambda_1\|\phi\|_1 + \lambda_2\|\theta\| + \lambda_3\|\phi\|
   \;+\;\gamma_0\|W\|_1 + \gamma_1\|W\| \nonumber\\
& +\;\mu_0\!\left\|\relu\!\left(+\tfrac{\partial f_\theta(\hat{h},\mathbf{a},m)}{\partial \hat{h}}\right)\right\|
   +\;\mu_1\!\left\|\relu\!\left(-\tfrac{\partial f_\theta(\hat{h},\mathbf{a},m)}{\partial m}\right)\right\| \nonumber\\
& +\;\mu_2\!\left\|\relu\!\left(-\tfrac{\partial g_\phi(\hat{h},\mathbf{a})}{\partial \hat{h}}\right)\right\|
   +\;\mu_3\,\big\|g_\phi(h_\star,\mathbf{a})-h_\star\big\|,
\label{eq:cx_gnn_loss}
\end{align}
where $m = \sum_{j\in\mathcal{N}_i}\hat{W}_{ij}\,g_\phi(\hat{h}_j,\mathbf{a}_j)$ is the aggregated synaptic input. The $\ell_1$ penalties $\lambda_0,\lambda_1$ on the MLP weights bias $f_\theta,g_\phi$ toward sparser, simpler functions; the $\ell_2$ penalties $\lambda_2,\lambda_3$ prevent weight explosion. The leak-monotonicity term $\mu_0$ penalises \emph{positive} slopes of $f_\theta$ in the postsynaptic state $\hat{h}$, encoding the prior that the drift contains a leak ($\partial f_\theta/\partial \hat{h} \le 0$); $\mu_1$ penalises negative slopes of $f_\theta$ in the aggregated input $m$, encoding the prior that postsynaptic drive increases with net excitatory input. Together, $\mu_0$ and $\mu_1$ force the discovered drift to be decreasing in self-voltage and increasing in synaptic input --- the two qualitative signatures of \cref{eq:cx_dynamics}. The term $\mu_2$ enforces the analogous monotonicity prior on $g_\phi$ in the presynaptic voltage. The normalisation term $\mu_3$ anchors $g_\phi$ at a reference subthreshold value $h_\star$, removing the multiplicative degeneracy between $g_\phi$ and $\hat{W}$ in \cref{eq:cx_gnn_dynamics}. Hyperparameters $\lambda_{\{0,1,2,3\}},\,\gamma_{\{0,1\}},\,\mu_{\{0,1,2,3\}}$ and the MLP widths/depths are tuned with the agentic optimisation procedure of \citep{allier_connectome_2026}.

\paragraph{Agentic hyperparameter optimisation.}\label{para:agent_gnn}
The GNN hyperparameter space (embedding dimension, MLP widths and depths, the $\lambda$, $\gamma$, $\mu$ coefficients, learning rates, and the curriculum on the rollout length $T$) is too large for grid search; we tune it with the closed-loop LLM-agent protocol of \citep{allier_connectome_2026}, using the deterministic-sweep rollout correlation $r_{\mathrm{roll},1k}$ (Pearson $r$ on the decoded heading at $\omega = 60^\circ/\mathrm{s}$, $T = 1000$ frames, warm-up $10$) as the primary score. The per-condition winning configurations selected by this procedure are listed in \cref{tab:cv_hyperparams} (Appendix B); empirical findings from this search --- including the embedding-dimension sweep --- are summarised in the Results.

\medskip
\noindent\emph{Why the agentic loop is methodologically primordial.} In high-dimensional recurrent systems, optimisation failure is not itself strong evidence against a mechanistic hypothesis. The non-convergence of a parameterisation may reflect poor hyperparameter selection, unstable training dynamics, or an unfavourable basin of the optimisation landscape rather than the absence of a valid solution. The agentic optimisation loop is therefore not merely a performance tool but a methodological requirement: without extensive closed-loop exploration of the training regime, it is easy to mistake optimisation failure for biological falsification.

\medskip
\noindent\emph{Two-stage agentic protocol.} We run each parameterisation through \emph{two} agentic loops with distinct objectives. The first loop is goal-driven: it searches the hyperparameter space until $r_{\mathrm{roll},1k}$ saturates at the behavioural ceiling, establishing that the parameterisation admits a converging recipe at all. The second loop is knowledge-driven: starting from the first loop's winner as a known floor, it conducts systematic single-axis sweeps across all hyperparameter families — learning rates, stabilisers, curriculum, regularisation, initialisation, architecture — under a strict causality rule (exactly one parameter changes per slot), interleaved with robustness batches (identical configuration, varied seeds) to separate genuine effects from seed variance. The two-stage split is essential because the first loop produces only the existence proof ("there is a recipe that converges") whereas the second produces the causal map the paper's interpretive claims rely on: which hyperparameters explain variance in the converged metric, which failure modes the parameterisation exhibits when training breaks, how robust the converged recipe is across seeds. Collapsing the two stages into a single optimisation pass yields a winning configuration without the mechanistic understanding such conclusions require.

\paragraph{Mutual-information estimators.}\label{method:MI}
We use mutual information to quantify how strongly each unit's
firing rate $r_i(t) = \sigma(\hat h_i(t))$ tracks the heading
$\theta(t)$ (or the angular velocity $\omega(t)$) over a held-out
rollout. For a single neuron we wrap $\theta$ to $(-\pi, \pi]$ via
$\theta \leftarrow \arg(e^{i\theta})$, discretise the wrapped angle
into $n_\theta$ equal-width angular bins and $r$ into $n_r$
equal-width bins, then compute the plug-in estimator in bits,
\begin{equation}
\hat I(r_i;\theta) = \sum_{b,c} \hat p(b,c)\,\log_2
\frac{\hat p(b,c)}{\hat p(b)\,\hat p(c)},
\label{eq:mi_plugin}
\end{equation}
where $\hat p$ is the empirical joint over the $(\theta\!\!-\!\!\mathrm{bin}, r\!\!-\!\!\mathrm{bin})$
histogram. For a group of $K$ neurons we use the cross-validated
lower bound
\begin{equation}
\hat I_{\mathrm{LB}}(\mathbf r_{1:K};\theta) =
H(\theta) - \tfrac{1}{|\mathcal D_{\mathrm{te}}|}
\sum_{(t,\theta)\in\mathcal D_{\mathrm{te}}}
-\log_2 q_\psi(\theta_b\mid \mathbf r(t)),
\label{eq:mi_joint_lb}
\end{equation}
where $q_\psi$ is an $L_2$-regularised multinomial logistic
regression with input $\mathbf r \in \mathbb R^K$ and output the
$n_\theta$ heading bins, trained on a stratified 5-fold split.
$\hat I_{\mathrm{LB}}$ is the population-MI lower bound used for
the per-cell-type and all-neurons joint MIs in
\cref{fig:hd_mi_summary,tab:hd_mi_summary,fig:omega_mi_summary,tab:omega_mi_summary}.
The same estimators apply to $\omega$ by replacing $\theta$ with
$\omega$ and using uniform binning over the 99\textsuperscript{th}
percentile of $|\omega|$ instead of circular wrapping.
The binning rule used throughout this paper is
$n_\theta = 64$ ($n_\omega = 64$) for the per-cell-type
fingerprints reported in
\cref{fig:hd_mi_summary,tab:hd_mi_summary,fig:omega_mi_summary,tab:omega_mi_summary},
and $n_\theta = 32$ ($n_\omega = 32$) for the per-neuron MI used
as the high/low-$\mi$ threshold in the R/L/D/Z partition of
\cref{method:four_classes,eq:four_classes}; the per-neuron split
is taken at coarser bin counts to keep the histogram estimator
well-populated at small $n_r$, while the joint $\hat I_{\mathrm{LB}}$
benefits from the finer $\theta$ resolution.

\paragraph{Functional classification of neurons (R/L/D/Z).}\label{method:four_classes}
Each trained network is partitioned into four functional classes by
combining three per-neuron scores measured on a $T = 10\,000$ frame
OU rollout: the heading mutual information $\mi_i = \hat I(r_i;
\theta)$ in bits using the plug-in estimator of \cref{eq:mi_plugin}
at $n_\theta = 32$ and $n_r = 20$; the readout magnitude $\omag_i =
\|\hat W^{\mathrm{out}}_{:,i}\|_2$, which is identically zero for
$i \ge n_{\mathrm{EPG}} = 46$ under the EPG-only readout; and the
angular-velocity modulation
\[
  \omod_i = \max\!\big(|\mathrm{corr}(r_i,\omega_+)|,\,
                       |\mathrm{corr}(r_i,\omega_-)|\big),
\]
where $\omega_{\pm} = \max(\pm\omega, 0)$ are the signed-half
angular velocities. Each axis is split at its sample median (the
$\omag$ median is taken over the non-zero $\omag$ entries to avoid
the trivial split that the EPG-only readout would otherwise force),
giving thresholds $(\mi^\star, \omag^\star, \omod^\star)$ that
assign every neuron to exactly one of
\begin{align}
  i \in \cls{R} &\iff \mi_i \ge \mi^\star \;\&\;\omag_i \ge \omag^\star
  \quad\text{(readout-relevant bump representation)} \nonumber\\
  i \in \cls{L} &\iff \mi_i \ge \mi^\star \;\&\;\omag_i < \omag^\star
  \quad\text{(latent HD encoders, invisible to the linear decoder)} \nonumber\\
  i \in \cls{D} &\iff \mi_i < \mi^\star \;\&\;\omod_i \ge \omod^\star
  \quad\text{(angular-velocity drivers)} \nonumber\\
  i \in \cls{Z} &\iff \mi_i < \mi^\star \;\&\;\omod_i < \omod^\star
  \quad\text{(uninvolved)}.
  \label{eq:four_classes}
\end{align}
The partition mirrors the zebrafish HD-circuit classification of
\citet{petrucco2023}, with swim modulation replaced by
angular-velocity modulation. Per-cell-type counts and the underlying
scatters / stacked-bar are reported in \cref{tab:cx_four_classes,fig:cx_four_classes}.


\section{Results}

\begin{figure}[H]
  \centering
  \includegraphics[width=\textwidth]{figure/fig_cx_anatomy_3d.png}
  \caption{\textbf{Anatomical morphology of the 156 hemibrain CX
  neurons modelled in this paper.} \textbf{(a)} 3-D skeleton render
  of the seven cell types (EPG, EPGt, PEN$_a$/PEN$_b$, $\Delta 7$,
  PEG, ER6; $n$ per type in the legend) overlaid on the neuropil
  outlines of EB, PB, FB, NO and BU$_{L/R}$ (grey silhouettes).
  View flipped so dorsal is up: PB at the top, FB in the middle,
  EB ring at the bottom, BU$_{L/R}$ on the right.
  \textbf{(b)} Same view restricted to the cell bodies. Each neuron
  contributes its soma 3-D mesh, recovered from the hemibrain
  segmentation by cropping the full per-neuron mesh to a sphere of
  radius $1.4\,r_{\mathrm{soma}}$
  around the soma node identified in the SWC skeleton (the soma is
  the only node with radius $\gg$ the segmental baseline). The two
  panels show \emph{spatially disjoint} compartments of the same set
  of neurons: in (a) the dense ring / moustache / fan structures
  (EB, PB, FB) are \emph{neuropils} --- regions of axo-dendritic
  arborisation with no cell bodies inside --- and in (b) the somata
  sit in the cortex around and above those neuropils, each connected
  to its neuropil arbor by a thin primary-neurite tract visible in
  (a). The EB ring in particular is therefore not a ring of cell
  bodies but the toroidal neuropil substrate the EPG / PEN /
  $\Delta 7$ arbors tile to support the ring-attractor dynamics.
  In (a), the EPG axons terminate in the PB (top, 18 glomeruli per
  the L1--L9 / R1--R9 numbering of \citet{hulse2021}) and their
  dendrites tile the 16 wedges of the EB ring (centre);
  $\Delta 7$ arborises across multiple PB glomeruli; ER6 lies
  largely outside the CX core (right). Cell bodies (b) sit in the
  cortex with no overlap with the neuropils in (a).
  Cortical cell bodies of each population sit outside the neuropil
  envelope in a stereotyped pattern --- e.g.\ the four ER6 somata
  cluster in one anatomically distinct location while the 42
  $\Delta 7$ somata wrap above the PB. The same colour assignment is used for
  the per-cell-type bars in
  \cref{fig:hd_mi_summary,fig:omega_mi_summary} and the row blocks
  in \cref{fig:connectome_summary}, so the morphology, connectivity
  and HD coding panels can be read against each other neuron type
  by neuron type.}
  \label{fig:cx_anatomy_3d}
\end{figure}

We compare four recurrent parameterisations on the same path-integration task, each with the EPG-only readout fixed and each evaluated under 10-fold cross-validation on independently seeded folds. The nominal training recipe runs without the EPG-ring TV regulariser ($\mathcal{L}_{\mathrm{tv}}$ of \cref{eq:cx_tv}; $\lambda_{\mathrm{tv}} = 0$); for the two convergent architectures (Known-ODE RNN and GNN) we additionally train paired $+$TV variants that turn the regulariser on, giving six conditions in total. The fully connected RNN and the frozen $W^{\mathrm{rec}}$ control are run only in the nominal no-TV configuration --- the former because TV did not improve its already-rare convergence, the latter because $\hat W^{\mathrm{rec}}$ is held fixed and the EPG-ring activity is determined by the encoder/decoder paths alone. The full set of per-condition learning rates and regularisation coefficients is listed in \cref{tab:cv_hyperparams} (Appendix B). Within each condition we report mean $\pm$ s.d.\ of $|r_{\mathrm{roll},1k}|$, path-integration accuracy $p_{\mathrm{acc}}$, EPG-bump FWHM, and constant-$\omega$ RMSE across the 10 folds, plus the number of folds that cross $|r_{\mathrm{roll},1k}| \ge 0.9$ (\cref{tab:cv_summary}).
\textbf{(i)~Known-ODE RNN} (\cref{method:RNN}, \cref{fig:evolution_pi}). Connectome topology and Dale signs are imposed via $\hat{W}^{\mathrm{rec}} = |\hat{S}|\odot\mathrm{sign}(W^{\mathrm{con}})$; the per-edge magnitudes $\hat{S}$, the four PEN velocity scalars, the encoder/decoder and biases train end-to-end. The nominal no-TV recipe converges 10/10 with $|r| = 1.000 \pm 0.000$ and RMSE $4.6 \pm 2.3^\circ$; the $+$TV variant also converges 10/10 ($|r| = 1.000 \pm 0.000$) but at a higher constant-$\omega$ RMSE of $13.3 \pm 7.3^\circ$. Adding TV therefore raises the constant-$\omega$ RMSE roughly three-fold without changing the Pearson-$r$ ceiling, but tightens the inter-fold $W^{\mathrm{rec}}$ cosine similarity from $0.66$ (no-TV) to $0.71$ ($+$TV) (\cref{fig:w_rec_comparison}).
\textbf{(ii)~Frozen-$W^{\mathrm{rec}}$ control} (\cref{fig:evolution_pi_frozen}). Same architecture with $\hat{W}^{\mathrm{rec}} \equiv W^{\mathrm{con}}$ held at the raw hemibrain values throughout training; only the encoder, decoder, biases, and PEN scalars train. 0/10 folds cross the $|r| \ge 0.9$ threshold; the mean $|r| = 0.669 \pm 0.094$ is non-trivial (the bump moves in the correct direction) but the path-integration accuracy collapses ($p_{\mathrm{acc}} = 0.29 \pm 0.13$, RMSE $= 95.7 \pm 16.3^\circ$): the connectome's nominal magnitudes set the right sign and topology of the recurrent operator but fail the gain that the heading readout requires.
\textbf{(iii)~Fully connected RNN} (\cref{fig:evolution_pi_fc}). The connectome topology mask is dropped; only the column-wise Dale sign of each presynaptic neuron is enforced. Only 3/10 folds converge ($|r| = 0.535 \pm 0.360$ across all folds, $p_{\mathrm{acc}} = 0.93 \pm 0.06$, RMSE $= 68.9 \pm 52.8^\circ$): without the connectome topology mask and with the EPG-only readout, the optimiser must discover both the support and the magnitudes from scratch and rarely succeeds.
\textbf{(iv)~Message-passing GNN} (\cref{fig:evolution_pi_gnn}). The dense recurrent operator is replaced by $f_\theta, g_\phi$ MLPs over a per-neuron embedding $\mathbf{a}_i$ (\cref{eq:cx_gnn_dynamics}); the leak and firing-rate non-linearity must emerge from these MLPs. The nominal no-TV recipe converges 9/10 ($|r| = 0.930 \pm 0.210$, one collapsed fold drags the mean), with the converged folds reaching RMSE $= 4.6 \pm 3.1^\circ$; the $+$TV variant converges 10/10 with $|r| = 1.000 \pm 0.001$ and RMSE $= 6.4 \pm 4.6^\circ$. The inter-fold $W$ similarity is $0.50$ no-TV and $0.52$ $+$TV, in both cases noticeably below the Known-ODE values.

\begin{table}[H]
  \centering
  \footnotesize
  \begin{tabular}{lccccc}
\toprule
Condition & $n_{\mathrm{conv}}$/10 & $|r_{\mathrm{roll},1k}|$ & $p_{\mathrm{acc}}$ & FWHM ($^\circ$) & RMSE ($^\circ$) \\
\midrule
Known-ODE RNN, no TV & 10/10 & $1.000\pm0.000$ & $0.996\pm0.002$ & $44.3\pm7.4$ & $4.6\pm2.3$ \\
Known-ODE RNN +TV & 10/10 & $1.000\pm0.000$ & $0.996\pm0.002$ & $84.0\pm8.0$ & $13.3\pm7.3$ \\
GNN, no TV & 9/10 & $0.930\pm0.210$ & $1.000\pm0.000$ & $41.3\pm12.9$ & $4.6\pm3.1$ \\
GNN +TV & 10/10 & $1.000\pm0.001$ & $0.999\pm0.001$ & $89.4\pm12.7$ & $6.4\pm4.6$ \\
Fully connected RNN & 3/10 & $0.535\pm0.360$ & $0.930\pm0.064$ & $31.6\pm9.2$ & $68.9\pm52.8$ \\
Frozen $W^{\mathrm{rec}}$ & 0/10 & $0.669\pm0.094$ & $0.294\pm0.134$ & $28.3\pm9.0$ & $95.7\pm16.3$ \\
\bottomrule
\end{tabular}

  \caption{\textbf{10-fold CV summary across the six trained
  conditions.} $n_{\mathrm{conv}}/10$ is the number of folds with
  $|r_{\mathrm{roll},1k}| \ge 0.9$ on a held-out constant-$\omega$
  rollout. The remaining columns are mean $\pm$ s.d.\ across the 10
  folds of the same condition. The same numbers drive every per-model
  paragraph below; the corresponding per-fold $W^{\mathrm{rec}}$
  analysis is given in \cref{fig:w_rec_comparison}.}
  \label{tab:cv_summary}
\end{table}

\begin{figure}[H]
  \centering
  \includegraphics[width=0.7\textwidth]{figure/fig_connectome_summary.png}
  \caption{\textbf{CX connectome $W^{\mathrm{con}}$: structure and
  per-cell-type edge-weight distributions.} The matrix is the
  effective connectome the trained Known-ODE / GNN models see,
  derived from the hemibrain export
  (\emph{papers/Code\_NN/Code\_NN/Data/Figure5/exported-traced-adjacencies-v1.2/},
  files \emph{traced-neurons.csv} and
  \emph{traced-total-connections.csv}); it follows the Hulse-2025
  preparation in which $\Delta 7$ and ER6 outputs are sign-flipped
  (inhibitory) and ER6 outputs are amplified $5\times$ before
  $W^{\mathrm{con}}$ is normalised to spectral radius $0.9$
  (\cref{eq:cx_w_rec}, \cref{method:RNN}).
  \textbf{(a)}~Signed $W^{\mathrm{con}}$, $z$-scored over non-zero
  entries and clipped to $\pm 3$; neuron-type blocks marked on
  both axes; red = excitatory, blue = inhibitory.
  \textbf{(b)}~Binary support $\mathrm{supp}(W^{\mathrm{con}})$
  (black = edge present, white = absent); the matrix is sparse
  (density 0.42, $10\,263$ non-zero edges out of $156^2 = 24\,336$).
  \textbf{(c)}~Per-cell-type \emph{incoming} edge-weight distributions
  (violin, mean marked): for each type $t$, the data are the non-zero
  $W^{\mathrm{con}}_{ij}$ with the post-synaptic neuron $i$ in $t$
  (variation across all pre-synaptic partners $j$).
  \textbf{(d)}~Per-cell-type \emph{outgoing} edge-weight distributions:
  for each type $t$, the data are the non-zero $W^{\mathrm{con}}_{ij}$
  with the pre-synaptic neuron $j$ in $t$ (variation across all
  post-synaptic partners $i$). Dale's law is visible in (d) as
  one-sided violins: every type is uniformly excitatory except
  $\Delta 7$ and ER6, whose outgoing distributions sit entirely
  below zero. Cell-type ordering and colour assignment match
  \cref{fig:hd_mi_summary,fig:omega_mi_summary}, so the rows there
  can be read off directly against the connectome statistics here.}
  \label{fig:connectome_summary}
\end{figure}

\begin{figure}[H]
  \centering
  \includegraphics[width=\textwidth]{figure/fig_evolution_drosophila_cx_pi_epg_no_tv_cv0.png}
  \caption{\textbf{Known-ODE RNN, no-TV recipe: connectome topology and Dale signs imposed, per-edge magnitudes learned.} Sign-locked RNN with $\hat{W}^{\mathrm{rec}} = |\hat{S}| \odot \mathrm{sign}(W^{\mathrm{con}})$, trained without the EPG-ring TV regulariser ($\lambda_{\mathrm{tv}} = 0$). Panel (c) is the same rollout-and-sort convention as the Known-ODE no-TV column of \cref{fig:kinographs_const} (column~b).
  \textbf{(a)}~GT connectome $W^{\mathrm{con}}$, $z$-scored over non-zero entries and clipped to $\pm 3$; neuron-type blocks marked on both axes.
  \textbf{(b)}~Learned $\hat{W}^{\mathrm{rec}}$ on the same scale.
  \textbf{(c)}~All 156 neurons over a constant-$\omega = +60^\circ/\mathrm{s}$ rollout, sorted by cell type; firing rates $z$-scored per neuron.
  \textbf{(d)}~PEN sub-population over the same rollout.
  \textbf{(e)}~EPG sub-population over the same rollout.
  \textbf{(f)}~Constant-$\omega$ rollout. Top: $\omega(t)$. Bottom: true heading (green) and decoded heading (black), wrapped. Title reports post-warmup Pearson $r$ and circular $\mathrm{RMSE}$.
  \textbf{(g)}~Same layout on a held-out OU test trial.
  \textbf{(h)}~Per-cell-type distribution of the subthreshold state $\hat{h}_i(t)$ (violin + mean $\pm 1$ SD), revealing where each population's activity sits relative to the sigmoid's linear regime.
  \textbf{(i)}~Integration gain: linear-fit slope of the unwrapped decoded heading vs true $\omega$, on a $15^\circ/\mathrm{s}$ grid over $[-180, +180]\setminus\{0\}$. Dashed line: $g = 1$; grey band marks the linear-tracking domain.
  }
  \label{fig:evolution_pi}
\end{figure}

\begin{figure}[H]
  \centering
  \includegraphics[width=\textwidth]{figure/fig_evolution_drosophila_cx_pi_epg_tv_cv0.png}
  \caption{\textbf{Known-ODE RNN, $+$TV companion to \cref{fig:evolution_pi}.} Same architecture, same fold (cv0), same panel layout, but trained with the EPG-ring TV regulariser ($\lambda_{\mathrm{tv}} = 10^{-2}$, \cref{eq:cx_tv}). The recurrent matrix in (b), the bump kinographs in (c)--(e), and the constant-$\omega$ / OU rollouts in (f, g) read the same as the no-TV companion; the constant-$\omega$ RMSE rises from $4.4^\circ$ to $14.9^\circ$ (panel (f) titles) and the inter-fold $W^{\mathrm{rec}}$ cosine cluster tightens from $0.66$ to $0.71$ (\cref{fig:w_rec_comparison}). Panel (c) corresponds to column~a of \cref{fig:kinographs_const}.}
  \label{fig:evolution_pi_tv}
\end{figure}

\begin{figure}[H]
  \centering
  \includegraphics[width=\textwidth]{figure/fig_evolution_drosophila_cx_pi_frozen_Wrec_epg_cv0.png}
  \caption{\textbf{Frozen-$W_{\mathrm{rec}}$ control.} Panels follow \cref{fig:evolution_pi}. The recurrent matrix is held at $\hat{W}^{\mathrm{rec}} \equiv W^{\mathrm{con}}$ throughout training; only $\hat{W}^{\mathrm{in}}$, $\hat{W}^{\mathrm{out}}$, biases, and the four PEN velocity scalars train. The decoded heading is flat (panel (f); $r = -0.37$, $\mathrm{RMSE} = 105^\circ$), the network does not respond to the OU velocity input (panel (g)), and the integration gain is zero across the full $\omega$ range (panel (i)). The model fails to form a bump, so the tuning, phase-shift, and FWHM panels (j-l) are omitted.
  }
  \label{fig:evolution_pi_frozen}
\end{figure}

\begin{figure}[H]
  \centering
  \includegraphics[width=\textwidth]{figure/fig_evolution_drosophila_cx_pi_fc_epg_cv0.png}
  \caption{\textbf{Fully connected RNN.} Panels follow \cref{fig:evolution_pi}. The connectome topology mask is dropped; $\hat{W}^{\mathrm{rec}} = |\hat{S}| \odot \mathrm{sign}\big(\sum_i W^{\mathrm{con}}_{i,:}\big)$ enforces only the dominant E/I sign per presynaptic neuron and the recurrent matrix is otherwise dense. The learned matrix in (b) carries weight outside the hemibrain edge support; heading accuracy is preserved ($r = 1.00$, $\mathrm{RMSE} = 12.4^\circ$). Panel (c) is the same rollout-and-sort convention as the fully connected column of \cref{fig:kinographs_const} (column~e).
  }
  \label{fig:evolution_pi_fc}
\end{figure}

\begin{figure}[H]
  \centering
  \includegraphics[width=\textwidth]{figure/fig_evolution_drosophila_cx_pi_gnn_epg_no_tv_cv0.png}
  \caption{\textbf{Message-passing GNN, no-TV recipe (\cref{method:GNN}).} Trained without the EPG-ring TV regulariser ($\lambda_{\mathrm{tv}} = 0$). Panel (c) is the same rollout-and-sort convention as the GNN no-TV column of \cref{fig:kinographs_const} (column~d). Panels
  \textbf{(a)--(i)} follow \cref{fig:evolution_pi}. The dense
  recurrent operator is replaced by
  $f_\theta\big(\hat{h}_i, \mathbf{a}_i, \sum_j \hat{W}_{ij}\,g_\phi(\hat{h}_j, \mathbf{a}_j)\big)$
  (\cref{eq:cx_gnn_dynamics}); the leak and firing-rate
  non-linearity emerge from $f_\theta, g_\phi$. Panel (b) shows the
  per-edge gain $\hat{W}_{ij}$ at the connectome edge positions (the
  GNN does not form a dense $\hat{W}^{\mathrm{rec}}$).
  \textbf{(j)}~Per-neuron latent embedding $\mathbf{a}_i \in \mathbb{R}^2$, coloured by cell type.
  \textbf{(k)}~Learned $f_\theta(\hat{h}_i, \mathbf{a}_i, m{=}0)$ for each neuron, plotted over its activity domain $[\mu_i - 2\sigma_i, \mu_i + 2\sigma_i]$ (snapshot from training). Per the GNN ODE (\cref{eq:cx_gnn_dynamics}), $f_\theta$ should approximate $-\hat{h}_i / \tau$ to recover the leak.
  \textbf{(l)}~Learned $g_\phi(\hat{h}_j, \mathbf{a}_j)$ over the same per-neuron domain (snapshot from training) --- the per-edge message that enters the GNN aggregation. The per-edge signed drive is $\hat{W}_{ij}\,g_\phi(\hat{h}_j,\mathbf{a}_j)$ (single-power, no squaring), so the effective synaptic sign is the product $\mathrm{sign}(\hat{W}_{ij}) \cdot \mathrm{sign}(g_\phi)$: Dale's law is no longer enforced per edge, though the dominant E/I row pattern is still encoded in $\hat{W}$. Because the squaring/positivity constraint is dropped, the $\pm$ sign degeneracy on $g_\phi$ is removed and the recovered firing-rate non-linearity is uniquely identified --- here sigmoidal in the standard ``depolarised-is-active'' convention, recovering the firing-rate non-linearity of \cref{eq:cx_dynamics}.
  Checkpoint shown is the epoch-5 best model of the EPG-readout
  tail-loss training.}
  \label{fig:evolution_pi_gnn}
\end{figure}

\begin{figure}[H]
  \centering
  \includegraphics[width=\textwidth]{figure/fig_evolution_drosophila_cx_pi_gnn_epg_tv_cv0.png}
  \caption{\textbf{Message-passing GNN, $+$TV companion to
  \cref{fig:evolution_pi_gnn}.} Same architecture and fold (cv0),
  trained with the EPG-ring TV regulariser
  ($\lambda_{\mathrm{tv}} = 5\!\times\!10^{-3}$, \cref{eq:cx_tv}).
  All panels follow \cref{fig:evolution_pi_gnn}. Panel (c)
  corresponds to column~c of \cref{fig:kinographs_const}. Visually
  the GNN $+$TV solution tightens the EPG kinograph (panel (e))
  along the EPG ring, slightly smooths the
  $(\hat h, g_\phi(\hat h))$ surface (panel (l)) toward the
  canonical sigmoid, and preserves the cell-type interleaving in
  the embedding (panel (j)); behavioural metrics ($|r| = 1.000$,
  RMSE rises from $4.7^\circ$ to $16.3^\circ$, panel (f)).}
  \label{fig:evolution_pi_gnn_tv}
\end{figure}

\paragraph{Known-ODE RNN: many circuits, same activity.}
The path-integration task fixes only the input--output map: angular
velocity $\omega(t)$ in, decoded heading $\hat \theta(t)$ out. The
Known-ODE RNN (\cref{method:RNN}, \cref{fig:evolution_pi}) imposes
the full ODE of \cref{eq:cx_dynamics} --- linear leak $-\hat h/\tau$,
sigmoid firing-rate non-linearity, hemibrain connectome topology +
Dale signs on $\hat W^{\mathrm{rec}}$ --- and trains only the
per-edge magnitudes $\hat S$. Across 10 independently initialised CV
folds (per-fold randomised simulation, training, and task seeds; the
winning recipe of \cref{para:agent_gnn}), every fold converges to
$r_{\mathrm{roll},1k} = 1.000 \pm 0.000$ at $T = 1000$ frames. The
recovered $\hat h$ trajectories cover both saturating tails of the
sigmoid (\cref{fig:function_dynamics}a, b: hexbin density of
$(\hat h, \sigma(\hat h))$ over a $T = 10{,}000$ frame OU rollout
shows non-zero occupancy out to $|\hat h| \approx 10$, with the bulk
of the per-neuron mass spread across the linear regime).

The 10 converged $\hat W^{\mathrm{rec}}$ matrices reach the same
behavioural ceiling but \emph{not} the same recurrent weights
(\cref{fig:w_rec_comparison}, leftmost column). The pairwise cosine
similarity over the connectome support (panel a) is $0.712 \pm
0.023$ between learned folds (45 off-diagonal pairs) and
$0.63 \pm 0.02$ against $W^{\mathrm{con}}$: the learned solutions
cluster more tightly with each other than with the connectome
template they sign-inherit from, but the inter-fold spread is still
substantial. Removing the TV regulariser (second column) leaves the
same pattern with a slightly lower inter-fold cosine ($0.663 \pm
0.023$); under the GNN the inter-fold cluster loosens further (third
and fourth columns: $0.518 \pm 0.024$ with TV, $0.502 \pm 0.021$
without). Panel b confirms
this at the per-cell-type-block level: some connection blocks are
tightly constrained across folds (median coefficient of variation
$\lesssim 0.5$, e.g.\ $\Delta 7 \to \mathrm{PEN}_b$,
$\mathrm{PEN}_b \to \mathrm{ER6}$), while others are essentially
free to wander ($\mathrm{CV} \gtrsim 1$, e.g.\ $\Delta 7 \to \Delta
7$, $\mathrm{EPGt} \to \mathrm{PEN}_{a/b}$). Panel c shows that
along the connectome support the trained matrices are systematically
$2$--$50\times$ larger than $W^{\mathrm{con}}$ (mean $\log_2$ ratio
positive for every block), and that this amplification pattern is
itself reproducible across folds ($\log_2$-ratio SDs $\lesssim
0.5$): training rescales the hemibrain template along its existing
edges rather than carving a different support.

Behaviour therefore does \emph{not} identify the per-edge weights:
the 10 folds reach the same $r_{\mathrm{roll},1k}$ via different
recurrent matrices that share an overall magnitude profile but
diverge at the per-edge level, and \cref{app:degeneracy_cx} shows
that per-cell-type partitions of the connectome support absorb
sum-zero perturbations at moderate amplitude
($\bar r = 0.99 \pm 0.05$ at $\lambda = 1$ across 10 models).
Many circuits, one activity. Tightening this equivalence class
further would require pinning the dynamics to more than just the
input--output map: per-neuron activity recordings across the
EPG / PEN / $\Delta 7$ populations would constrain the voltage
trajectories themselves and collapse the per-edge ambiguity that
the heading readout alone cannot resolve.

\begin{figure}[H]
  \centering
  \includegraphics[width=\textwidth]{figure/fig_w_rec_comparison.png}
  \caption{\textbf{Comparison of the trained
  $\hat W^{\mathrm{rec}}$ to each other and to the connectome
  template $W^{\mathrm{con}}$, across four parameterisations $\times$
  10 CV folds.}
  Columns: Known-ODE no-TV, Known-ODE $+$TV, GNN no-TV, GNN $+$TV.
  \textbf{(a)} Per-condition $11 \times 11$ pairwise cosine similarity
  over the connectome support, ordered $W^{\mathrm{con}}$ first, then
  the 10 CV folds. Diagonal $= 1.0$. Inter-fold mean off-diagonal
  (bottom-right $10 \times 10$ block): $0.663 \pm 0.023$ (Known-ODE
  no-TV), $0.712 \pm 0.023$ (Known-ODE $+$TV), $0.502 \pm 0.021$
  (GNN no-TV), $0.518 \pm 0.024$ (GNN $+$TV). First row / column
  off-diagonal ($\hat W_k$ vs.\ $W^{\mathrm{con}}$) is $0.63 \pm 0.02$
  for Known-ODE $+$TV and similar across columns. In every condition
  learned solutions cluster more tightly with each other than with the
  raw connectome template; the cluster is tightest for the Known-ODE
  RNN and loosens as the dynamical prior is relaxed (GNN).
  \textbf{(b)} Per-block median coefficient of variation across the
  10 folds, $\sigma_k |\hat W_k| / \langle |\hat W_k| \rangle_k$,
  for each (post-type, pre-type) cell-type pair on the same $7
  \times 7$ block layout as \cref{fig:connectome_summary}c, d. Low
  CV $=$ tightly constrained block; high CV $=$ free block. Empty
  grey cells have no connectome edges between that cell-type pair.
  \textbf{(c)} Per-block $\log_2$ ratio of the mean learned magnitude
  to the mean connectome magnitude,
  $\log_2 \big[\langle |\hat W_k|\rangle / \langle |W^{\mathrm{con}}|\rangle\big]$,
  annotated as $\mu \pm$ SD across the 10 folds. All blocks are
  uniformly positive (trained solutions amplify the connectome along
  its existing support), with the amplification factor varying by an
  order of magnitude across cell-type pairs and tightly conserved
  across folds within each parameterisation.}
  \label{fig:w_rec_comparison}
\end{figure}

\begin{figure}[H]
  \centering
  \includegraphics[width=\textwidth]{figure/fig_function_dynamics_combined_epg.png}
  \caption{\textbf{Operating range of the leak/drift and the
  firing-rate non-linearity along a natural OU velocity rollout
  ($T = 10\,000$ frames) for the six parameterisations under the
  EPG-only readout.} Hexbin density (log10 count, Blues) of the
  per-(neuron, time) samples; red curve is the static function over
  the plotted $\hat h$ range (with $\pm 1$~SD across the 156 neurons
  where applicable). Layout: 3 rows $\times$ 4 columns. Each row
  pairs two conditions side by side, with each condition spanning
  two columns (drift on the left, per-neuron firing-rate
  non-linearity on the right). Row 1: Known-ODE no-TV (left pair),
  Known-ODE $+$TV (right pair). Row 2: GNN no-TV (left pair),
  GNN $+$TV (right pair). Row 3: fully connected (left pair),
  frozen $W^{\mathrm{rec}}$ (right pair).
  \textbf{Row 1 (Known-ODE $\pm$TV).} Hardcoded linear leak and
  sigmoid; the rollout sweeps the full $|\hat h| \lesssim 10$ domain
  --- the bump traverses every glomerular phase as heading changes,
  painting the linear leak and saturating $\sigma$ in both tails.
  \textbf{Row 2 (GNN $\pm$TV).} The learned drift $f_\theta$
  deviates from $-\hat h/\tau$ over the visited $\hat h$ range and
  the learned signalling $g_\phi$ deviates from the canonical
  sigmoid except at a coarse qualitative level enforced by the
  $\partial g_\phi / \partial v \ge 0$ monotonicity prior. The
  operating range concentrates in a narrow $\hat h$ band near a
  flat plateau; the steep non-linear regimes are sparsely probed
  by the bump-encoding sub-population.
  \textbf{Row 3 (fully connected and frozen $W^{\mathrm{rec}}$).}
  Fully connected uses the same hardcoded $-\hat h/\tau$ and
  $\sigma$ as Known-ODE, but the operating range visited is shaped
  by a different learned $\hat W^{\mathrm{rec}}$ that populates
  off-block entries absent from the hemibrain wiring. The frozen
  $W^{\mathrm{rec}}$ control fails to produce a bump, and the
  hexbin collapses onto a tight band around the resting state.}
  \label{fig:function_dynamics}
\end{figure}

To probe how HD information is parcellated across cell types, we
estimate the mutual information $I(\sigma(\hat h_i); \theta)$
between each neuron's rate and heading along a $T = 10\,000$ frame
OU rollout ($n_\theta = 64$ angular bins for the per-cell-type
fingerprint of \cref{fig:hd_mi_summary,tab:hd_mi_summary};
$n_\theta = 32$ for the per-neuron MI thresholding of the R/L/D/Z
partition, \cref{method:MI,method:four_classes}). On the cv0 fold
of the Known-ODE $+$TV recipe (\cref{tab:hd_mi_summary})
the whole-population joint MI is $I_{\mathrm{pop}} \approx 3.2$
bits, predicting a Gaussian-channel per-step posterior
$\sigma_{\theta\mid\hat h} \approx \sigma_\theta\,2^{-I_{\mathrm{pop}}}
\approx 11^\circ$ --- in the same band as the directly measured
circular RMSE, confirming the MI estimator tracks the network's
actual heading precision. The per-neuron champion is EPGt
($\bar I = 2.33$ bits) followed by EPG and $\Delta 7$
($\bar I \approx 2.0$ bits each); ER6 is HD-tuned at $\bar I = 1.70$
bits and is therefore \emph{not} silent in this recipe. Per-cell-type
cloning is summarised by the redundancy factor
$\rho_t = \sum_{i \in t} I_i / I_{\mathrm{joint},t}$
(\cref{tab:hd_mi_summary}); large $\rho$ means same-type cells carry
near-identical $\theta$-tuning, so per-neuron MIs sum sub-additively
into the joint. The EPG block has $\rho \approx 39$:
the 46 EPG cells contribute only $\sim 46/\rho \approx 1$
effective independent encoder, because heavy cloning within each
EB wedge saturates the joint against
$H(\theta) = \log_2 64 = 6$ bits at $I_{\mathrm{joint},\mathrm{EPG}}
= 2.61$ bits. This also explains why the all-neurons joint MI
($\sim 3.2$ bits) is so much smaller than $\sum_{t,i} I_i$
($\approx 280$ bits naively): the code is heavily cloned and
capped near the angular-entropy ceiling rather than scaling with
neuron count.
The remaining three trained parameterisations
(\cref{fig:evolution_pi_frozen},\cref{fig:evolution_pi_fc},\cref{fig:evolution_pi_gnn})
expose three further layers of this degeneracy, addressed in turn
below.

\paragraph{Frozen-$W^{\mathrm{rec}}$ control: the raw connectome supplies topology, training supplies the gain.}
A natural reading of the Known-ODE RNN's success is that the
connectome itself --- not the trained magnitudes --- carries the
path-integration computation. The frozen-$W^{\mathrm{rec}}$ control
(\cref{method:RNN} with $\hat W^{\mathrm{rec}} \equiv W^{\mathrm{con}}$
held at the raw hemibrain values throughout training,
\cref{fig:evolution_pi_frozen}) falsifies that reading on the
behavioural metric: 0/10 CV folds cross $|r_{\mathrm{roll},1k}| \ge
0.9$ and the path-integration accuracy collapses
($p_{\mathrm{acc}} = 0.29 \pm 0.13$, constant-$\omega$ RMSE
$= 95.7 \pm 16.3^\circ$, integration gain $\approx 0$,
\cref{tab:cv_summary}). The non-trivial residual mean
$|r| = 0.669 \pm 0.094$ has a clean mechanistic reading rather
than being noise: the decoded heading tracks the true heading at
the right \emph{rotation direction} but the wrong overall gain,
so the bump moves through the right sequence of EPG glomeruli at
the wrong rate. The raw hemibrain weights therefore fix the
\emph{sign} and \emph{topology} of the recurrent operator (enough
for an $\omega$-modulated rotation in the cos/sin readout) while
the trained per-edge magnitudes of the Known-ODE RNN supply the
\emph{gain} that makes the rotation a 1:1 path integrator. This
is the complementary statement to \citet{hulse2026}: their
hidden-symmetry classes assume an attractor exists and ask which
microscopic perturbations preserve it; we show that training, not
the topology alone, brings the attractor into existence.

\paragraph{Fully connected RNN: removing the topology mask collapses convergence under the EPG-only readout.}
Replacing the topology mask with a column-wise Dale sign constraint
($\hat W^{\mathrm{rec}} = |\hat S| \odot
\mathrm{sign}\big(\sum_i W^{\mathrm{con}}_{i,:}\big)$,
\cref{fig:evolution_pi_fc}) allows the recurrent matrix to populate
every off-block entry absent from the hemibrain wiring. Across 10
CV folds, only 3/10 converge to $|r_{\mathrm{roll},1k}| \ge 0.9$
(\cref{tab:cv_summary}; mean $|r| = 0.535 \pm 0.360$,
$p_{\mathrm{acc}} = 0.93 \pm 0.06$, RMSE $= 68.9 \pm 52.8^\circ$).
The fully connected architecture is therefore rejected on
\emph{convergence reliability}, not on solution quality given
convergence: of the three folds that do converge, the trained
$\hat W^{\mathrm{rec}}$ populates off-block entries absent from
the hemibrain wiring (\cref{fig:evolution_pi_fc}b) and the
ring-attractor bump migration is preserved
(\cref{fig:kinographs}); the remaining seven folds fail to find
any heading-tracking solution under the EPG-only readout. The
connectome's pattern of missing edges is therefore not strictly
necessary for the computation when it can be found, but with the
EPG-only readout the optimiser must discover both the support and
the magnitudes from scratch and rarely succeeds under matched
training budgets. On the cv0 fold (which does converge), the
per-cell-type HD MI (\cref{fig:hd_mi_summary}) reaches
$I_{\mathrm{pop}} = 3.20$ bits, but the internal allocation
diverges from the Known-ODE solution: EPG falls to $\bar I = 0.61$
bits and the per-neuron champions become $\Delta 7$ ($\bar I =
1.13$ bits) and ER6 ($\bar I = 1.19$ bits) --- the same
``HD information reroutes onto the inhibitory backbone'' pattern
visible in the Hulse hidden-symmetry analysis.

\paragraph{GNN with learned neuronal dynamics: the implementation is not uniquely identified by the task.}
The GNN (\cref{method:GNN}, \cref{fig:evolution_pi_gnn}) keeps the
connectome topology mask but replaces the hardcoded leak and
sigmoid of the Known-ODE RNN with learned functions
$f_\theta, g_\phi$. Across 10 CV folds the nominal no-TV recipe
converges 9/10 to $|r_{\mathrm{roll},1k}| = 0.930 \pm 0.210$
(one collapsed fold accounts for the non-zero spread,
\cref{tab:cv_summary}); turning the TV regulariser on gives the
$+$TV variant which converges 10/10 to $|r| = 1.000 \pm 0.001$
with RMSE $= 6.4 \pm 4.6^\circ$. Despite the much weaker dynamical prior, the
trained GNN therefore reproduces the same ring-attractor behaviour
as the Known-ODE RNN, and does so with a noticeably more linear
integration gain: in \cref{fig:evolution_pi_gnn}i the decoded
heading rate tracks the true $\omega$ along the $y = x$ line over
the full $\pm 180^\circ/\mathrm{s}$ tested range (slope $\approx 1$,
small scatter), whereas the Known-ODE RNN
(\cref{fig:evolution_pi}i) shows visible saturation / sigmoidal
distortion at large $|\omega|$. This is the expected consequence of
the GNN dropping the hard-coded sigmoid: with $f_\theta, g_\phi$
free to be more linear in the bump-encoding regime, the network
can match the per-step heading update over a wider angular-velocity
range without paying the saturating-tail penalty that the fixed
$\sigma$ imposes on the Known-ODE solution. The recovered dynamics
differ substantially from the reference Known-ODE model. The learned drift
$f_\theta$ (\cref{fig:evolution_pi_gnn}k) is not the linear leak
$-h/\tau$, and the learned signalling function $g_\phi$
(\cref{fig:evolution_pi_gnn}l) does not reproduce the canonical
sigmoid except at a coarse qualitative level enforced by the
monotonicity prior. Yet the resulting network still performs path
integration with near-identical behavioural accuracy. Moreover, the rollout probes only a narrow operating region of the
learned functions (\cref{fig:function_dynamics}g, h): most neurons
remain near a low-activity plateau while only a small
bump-encoding sub-population traverses the strongly nonlinear
regime. Large regions of the learned dynamical surfaces are
therefore effectively unconstrained by the task. The GNN thus exposes a deeper layer of degeneracy than the
per-weight or connectome-zero ambiguities of the previous models.
The path-integration task does not uniquely identify the
neuronal dynamics themselves: multiple drift and signalling
functions can realise the same ring-attractor computation.
The same degeneracy appears in the learned per-neuron embedding
$\mathbf{a}_i \in \mathbb{R}^4$ (\cref{fig:evolution_pi_gnn}j
shows its first two coordinates): the embedding does \emph{not}
cluster into separate islands per cell type. Quantitatively, the
silhouette score of $\mathbf{a}_i$ with cell-type labels is
$s = -0.04$ on the full 4D embedding (and $s < 0$ on the 2D
projection of panel j), where $s = 0$ marks no separation and
$s < 0$ marks anti-clustering --- a single anatomical label is on
average \emph{further} from its own type centroid than from the
nearest other-type centroid. The path-integration objective
therefore does not force the network to specialise its dynamics
per cell type: the behaviourally relevant axes of $\mathbf{a}_i$
correspond to a finer per-neuron functional role (e.g.\ position
on the EPG ring) rather than to the seven anatomical labels,
leaving the cell-type allocation of HD information as an emergent
property of the trained $\hat W^{\mathrm{rec}}$ rather than an
explicit prior of the GNN.
The per-cell-type HD code (\cref{fig:hd_mi_summary}) shows that
the GNN $+$TV solution carries HD information across the same
cell-type backbone as the Known-ODE RNN --- EPG, EPGt, PEN, $\Delta
7$ and PEG all attain per-neuron MI $> 0.4$ bits --- but with a
lower overall magnitude per neuron than the Known-ODE solution
(EPG $0.57$ vs.\ $2.20$ bits; $\Delta 7$ $0.94$ vs.\ $2.00$ bits).
The exception is ER6, which the GNN $+$TV essentially silences
($\bar I = 0.22$ bits) while the Known-ODE RNN keeps it
HD-tuned ($\bar I = 1.70$ bits). The GNN's lower per-neuron MI
budget is consistent with its concentrated operating range
(\cref{fig:function_dynamics}): most cells stay near the flat
plateau of the learned $g_\phi$ and contribute little to the
heading code, while a small bump-encoding sub-population in the EPG
ring carries the signal across to the readout.

\paragraph{Identifiable computation, non-identifiable implementation.}
Three of the four parameterisations (Known-ODE RNN, fully connected
RNN, GNN) reach essentially identical heading accuracy on path
integration but realise the computation through qualitatively
different machinery: different per-edge weights
(\cref{app:degeneracy_cx}), different presence or absence of
off-block edges (Known-ODE vs.\ fully connected), and different
per-neuron intrinsic functions (RNN's hardcoded $-h/\tau,\sigma$
vs.\ GNN's learned $f_\theta, g_\phi$). The fourth, frozen-$W^{\mathrm{rec}}$
control fails outright (\cref{fig:evolution_pi_frozen}), ruling out
the simplest reading that the raw connectome alone instantiates the
algorithm. What remains \emph{identifiable} across the three
that converge?
\textbf{(i)} The functional architecture: all three implement a
ring-attractor with a bump that propagates as heading changes,
visible as a diagonal bump migration in \cref{fig:kinographs}.
\textbf{(ii)} The input--output routing: in all three trained
networks, the EPG / PEN / $\Delta 7$ neurons participate in the
bump while ER6 / PEG cells do not --- a consequence of the
$\hat W^{\mathrm{in}}$ row structure they share (stimulus routed
to PEN sub-populations only, initial-heading impulse free for all).
What is \emph{not} identifiable: per-edge $\hat W^{\mathrm{rec}}$
(\cref{app:degeneracy_cx}), the form of the per-neuron drift, the
form of the firing-rate non-linearity, and the anatomical-neuron
$\to$ preferred-phase mapping (different trained networks assign
different anatomical neurons to the same phase $\varphi$ along the
ring). Optimising the path-integration input--output map fixes
the ring-attractor computation; the implementation remains
under-determined at every level we probe.

\paragraph{Cell-type allocation of HD information.}
The four trainable parameterisations all reach high heading
precision ($I_{\mathrm{pop}}$ between $2.9$ and $3.6$ bits;
\cref{fig:hd_mi_summary}) while the frozen-$W^{\mathrm{rec}}$
control collapses to $I_{\mathrm{pop}} \approx 0.08$ bits.
Where the four trained solutions differ is in how the per-neuron
MI budget is distributed across cell types.
\textbf{(i)}~The Known-ODE RNN $+$TV places strong HD information
on EPG, EPGt, $\Delta 7$ and PEG (per-neuron mean $\bar I = 2.20,
2.33, 2.00, 1.92$ bits) with PEN sub-populations carrying
intermediate values ($\bar I = 1.32 / 1.00$ bits); this is
consistent with the canonical CX picture
\citep{seelig2015,kim2017,green2017} in which EPG, EPGt, $\Delta
7$ and PEG carry HD content. ER6, however, is here HD-tuned at
$\bar I = 1.70$ bits --- a quantitative deviation from the
canonical picture in which ER6 is tonic / not HD-tuned. The
no-TV recipe inherits the same ER6 anomaly
(\textbf{(ii)} below, $\bar I = 1.74$ bits).
\textbf{(ii)}~Removing the TV regulariser leaves the same
qualitative ranking with slightly lower per-neuron values
(Known-ODE no-TV: EPG $1.80$, $\Delta 7$ $1.95$, ER6 $1.74$ bits).
\textbf{(iii)}~The fully connected RNN shifts the per-neuron
champion to $\Delta 7$ ($1.13$ bits) and ER6 ($1.19$ bits), with
EPG and PEN carrying $\le 0.6$ bits --- a reallocation away from
the EPG-centred canonical assignment despite the same EPG-only
readout. \textbf{(iv)}~The GNN $+$TV keeps the canonical
EPG / $\Delta 7$ / PEG / EPGt backbone but at a uniformly lower
per-neuron budget ($\bar I \in [0.5, 1.2]$ bits) and silences ER6
($\bar I = 0.22$ bits); on this single cell type the GNN $+$TV is
in fact closer to the canonical tonic-ER6 picture than either
Known-ODE recipe is. Removing TV from the GNN restores ER6
($\bar I = 1.63$ bits) and pushes the joint MI to its highest value
in the sweep ($I_{\mathrm{pop}} = 3.64$ bits). The biological
division of labour is therefore \emph{not} forced by the
path-integration task alone: the same input--output map admits
trained solutions that reach near-identical $I_{\mathrm{pop}}$
through measurably different per-cell-type allocations, with the
Known-ODE $+$TV configuration the closest to the canonical
EPG / $\Delta 7$ partition reported by calcium imaging.


\paragraph{Relation to Hulse et~al.\ 2026.}
The close companion to our analysis is the
\emph{hidden-symmetries} framework of \citet{hulse2026}, which
shows how ring-attractor dynamics can be preserved under
heterogeneous connectivity through hidden symmetries embedded in
cloned network architectures. Their study combines analytically
constructed symmetry-preserving networks with unconstrained
optimisation results, arguing that the CX attractor is robust to
substantial microscopic heterogeneity. Our analysis asks a
complementary question: given successful behavioural recovery on
the same path-integration task, which aspects of the recovered
implementation are actually identifiable? Our
per-$(i,\alpha,\mathrm{instance})$ null directions
(\cref{app:degeneracy_cx}) overlap closely with the clone structure
highlighted in their framework, but we further show that the
degeneracy extends beyond connectivity to the intrinsic neuronal
dynamics themselves.

\textbf{(i)~Per-cell-type perturbations are locally null but do not
remain invariant under full collapse.}
Mixing clones across PB glomeruli (the per-$(i,\alpha)$ partition)
breaks the readout at $r = 0.36$. Per-cell-type groupings are not
identified as clones by \citet{hulse2026} --- their analysis works
at the per-glomerulus instance level --- and our empirical test
confirms it: lumping same-type neurons together is too coarse to
saturate the null space.
\textbf{(ii)~Per-unit clones more closely track the invariant
subspace.} Our
per-$(i,\alpha,\mathrm{instance})$ partition (\cref{app:degeneracy_cx})
groups the same-type, same-PB-glomerulus neurons that Hulse
et~al.\ identify as clones, and sum-preserving collapses inside
each group leave the rollout at $r = 0.83$ even after zeroing
$50.2\%$ of the edges --- the structural symmetry is operationally
real. The natural unit of redundancy is the glomerulus, not the
cell type, exactly as the symmetry framework predicts.
\textbf{(iii)~The connectome's wiring is necessary
but not sufficient.} The frozen-$W^{\mathrm{rec}}$ control
(\cref{fig:evolution_pi_frozen}) fails to integrate, so the
symmetry classes of the raw connectome \emph{cannot} on their own
instantiate the attractor: learned per-edge magnitudes inside those
classes are doing real work. This is the complementary statement
to \citet{hulse2026}: their robustness proof assumes a working
attractor; we show that training, not topology, supplies it.
\textbf{(iv)~The functional form is also non-unique.} The
hidden-symmetries analysis fixes a sigmoidal $\sigma$ and a linear
leak as part of the model; our GNN result
(\cref{fig:evolution_pi_gnn}) shows that behaviourally indistinguishable
dynamics arise from a learned $(f_\theta, g_\phi)$ pair that
\emph{neither} reduces to $-h/\tau$ nor to the canonical sigmoid.
The symmetry classes therefore explain one layer of degeneracy ---
hidden symmetries within a fixed neuronal model --- whereas our GNN
result adds a second layer: the neuronal dynamics themselves are
also non-identifiable. The per-cell-type HD MI analysis
(\cref{fig:hd_mi_summary}, \cref{tab:hd_mi_summary}) sharpens this
point: alternative signalling functions $(f_\theta, g_\phi)$ do
recover the behaviour, but they do so with a measurably different
per-cell-type allocation of HD information --- e.g.\ the GNN $+$TV
silences ER6 and the fully connected RNN reallocates HD onto
$\Delta 7$ and ER6 at the expense of EPG --- while calcium imaging
\citep{seelig2015,kim2017,green2017} reports EPG, EPGt, PEN, PEG
and $\Delta 7$ as HD-tuned.


\begin{figure}[H]
  \centering
  \includegraphics[width=\textwidth]{figure/fig_traces_const_omega_epg.png}
  \caption{\textbf{Voltage traces of the six parameterisations under a
  deterministic constant angular velocity stimulus $\omega = 60^\circ$/s
  ($T = 1\,500$ frames at $dt = 0.01$ s $= 15$ s $= 2.5$ full turns).}
  Same per-model trace layout as \cref{fig:task_traces}: 5
  representative neurons per cell type stacked vertically, with the
  L / R PEN stim overlay (light red / light blue) on the first
  PEN$_a$ / PEN$_b$ row of each cluster; bottom strip shows wrapped
  HD GT in green vs.\ decoded in black. Under a constant-$\omega$
  drive the head direction sweeps linearly and the bump migration is
  periodic, so each EPG / PEN / $\Delta 7$ neuron rings up and down at
  a fixed phase offset --- the canonical ring-attractor signature.
  3$\times$2 panel: \textbf{(a)} Known-ODE $+$TV,
  \textbf{(b)} Known-ODE no-TV, \textbf{(c)} GNN $+$TV,
  \textbf{(d)} GNN no-TV, \textbf{(e)} fully connected,
  \textbf{(f)} frozen $W^{\mathrm{rec}}$. The five converged
  parameterisations reproduce the constant-rate HD sweep with
  periodic bump migrations; the frozen-$W^{\mathrm{rec}}$ panel shows
  no bump and a decoded HD that drifts off the linear ground truth.}
  \label{fig:task_traces_const}
\end{figure}

\begin{figure}[H]
  \centering
  \includegraphics[width=\textwidth]{figure/fig_traces_combined_epg.png}
  \caption{\textbf{Voltage traces of the six parameterisations under
  the natural OU velocity stimulus.} Each tile is the same task-traces
  layout (5 representative neurons per cell type stacked vertically;
  PEN$_a$ / PEN$_b$ clusters carry the L / R stim overlay in light red
  / blue; bottom strip shows wrapped HD GT in green, decoded in
  black). 3$\times$2 panel: \textbf{(a)} Known-ODE $+$TV,
  \textbf{(b)} Known-ODE no-TV, \textbf{(c)} GNN $+$TV,
  \textbf{(d)} GNN no-TV, \textbf{(e)} fully connected,
  \textbf{(f)} frozen $W^{\mathrm{rec}}$. The five converged
  parameterisations reproduce path integration on this rollout
  (decoded HD tracks GT) but the underlying subthreshold dynamics
  differ visibly: the four RNN variants drive every neuron through a
  full bump-traversal cycle as heading changes, while the GNN
  concentrates most cells in a near-zero silent plateau and lets only
  a small bump-encoding sub-population carry the integration
  (cf.~\cref{fig:function_dynamics}).}
  \label{fig:task_traces}
\end{figure}

\begin{figure}[H]
  \centering
  \includegraphics[width=0.85\textwidth]{figure/fig_kinographs_const_omega_epg.png}
  \caption{\textbf{All-neurons-by-time kinographs of the six
  parameterisations' subthreshold state under a deterministic constant
  angular velocity stimulus $\omega = 60^\circ$/s ($T = 1\,500$ frames
  at $dt = 0.01$ s $= 15$ s $= 2.5$ full turns), in three row
  orderings.} Each tile shows the per-neuron z-scored
  $(\hat h_i(t) - \langle\hat h_i\rangle)/\mathrm{std}(\hat h_i)$
  for all 156 CX neurons; colour fixed at $\pm 3$ z-scores so the
  columns are directly comparable regardless of each model's
  absolute voltage scale. Columns (left to right): Known-ODE no-TV,
  Known-ODE $+$TV, GNN no-TV, GNN $+$TV, fully connected, frozen
  $W^{\mathrm{rec}}$.
  \textbf{Row 1 (cell-type sort, ER6 $\to$ EPG top-to-bottom).}
  Same row ordering as \cref{fig:task_traces}. Known-ODE and the
  fully connected RNN show a bump migrating through the EPG / PEN /
  $\Delta 7$ blocks as heading changes; the GNN retains the
  bump-encoding sub-population but the cell-type ordering scrambles
  the bump migration into apparent noise, because same-type cells
  encode different glomeruli. The frozen $W^{\mathrm{rec}}$ column
  shows no organised bump under any sort, consistent with the failed
  rollout.
  \textbf{Row 2 (cell-type primary, preferred-phase secondary
  within type).} Cell-type blocks preserved (ER6 $\to$ EPG), but
  rows within each type are sorted by their preferred heading phase
  $\varphi_i = \arg\!\big(\sum_t (\hat h_i(t) - \langle \hat h_i
  \rangle)\,e^{i\theta(t)}\big)$. Each block now reads as a small
  ring: the per-block diagonals all share the same instantaneous
  phase, making the ring-attractor bump visible without losing the
  anatomical type labels.
  \textbf{Row 3 (pure preferred-phase sort, low-$\varphi$ at top).}
  Under this ordering the five converged columns reveal the same
  diagonal bump-migration signature --- the canonical ring-attractor
  readout. The constant-$\omega$ drive fixes the bump-migration slope
  to $\omega$, so the phase-sorted diagonals are clean and parallel.
  Each converged parameterisation resolves the task with a
  ring-attractor structure built on a \emph{different} mapping from
  anatomical neuron-index to $\varphi$.}
  \label{fig:kinographs_const}
\end{figure}

\begin{figure}[H]
  \centering
  \includegraphics[width=0.85\textwidth]{figure/fig_kinographs_both_sorts_epg.png}
  \caption{\textbf{Same six-column layout and row orderings as
  \cref{fig:kinographs_const}, but under the natural OU velocity
  stimulus ($T = 10\,000$ frames $\approx 100$ s).} Replacing the
  deterministic constant-$\omega$ drive of \cref{fig:kinographs_const}
  with the stochastic OU drive used for training preserves the same
  ring-attractor signature in all five converged columns. The
  phase-sorted diagonals (bottom row) now wander with the OU
  heading trajectory and pause whenever the stimulus inserts a
  standing interval; the frozen-$W^{\mathrm{rec}}$ column remains
  unstructured under every sort.}
  \label{fig:kinographs}
\end{figure}


\begin{figure}[H]
  \centering
  \includegraphics[width=0.7\textwidth]{figure/fig_hd_mi_summary_nbins64.png}
  \caption{\textbf{Per-cell-type head-direction mutual information
  across the six parameterisations}, evaluated along the natural OU
  velocity rollout ($T = 10\,000$ frames, $n_\theta = 64$ angular
  bins). One row per model. \textbf{Left (a, c, e)} per-neuron MI
  $I(\sigma(\hat h_i);\theta)$ in bits (bars = per-cell-type mean,
  dots = individual neurons); \textbf{right (b, d, f)} joint per-type
  MI $I(\{\sigma(\hat h_i)\}_{i\in t};\theta)$ (CV-logreg lower bound),
  with the population joint across all 156 neurons annotated inline.
  Cell types are ordered by the Known-ODE per-neuron mean so vertical
  comparison reads as a per-cell-type fingerprint. Rows correspond
  to the six parameterisations (Known-ODE $\pm$TV, GNN $\pm$TV,
  fully connected, frozen $W^{\mathrm{rec}}$). The five trainable
  parameterisations reach population joint MI between $2.9$ and
  $3.6$ bits, predicting a Gaussian-channel per-step posterior
  $\sigma_{\theta\mid\hat h} = \sigma_\theta\,2^{-I_{\mathrm{pop}}}
  \sim 10\text{--}12^\circ$ ($\sigma_\theta = \pi/\sqrt 3
  \approx 1.81$~rad), of the same order as the directly measured
  circular RMSEs; the frozen $W^{\mathrm{rec}}$ control collapses
  to $I_{\mathrm{pop}} \approx 0.08$ bits. The five trained
  solutions distribute HD information across cell types differently
  (Known-ODE $+$TV preserves the canonical EPG / $\Delta 7$ /
  EPGt / PEG pattern; the GNN $+$TV silences ER6; the fully
  connected RNN reallocates HD onto $\Delta 7$ and ER6 at the
  expense of EPG; cf.\ \cref{tab:hd_mi_summary}).}
  \label{fig:hd_mi_summary}
\end{figure}

\begin{table}[H]
  \centering
  \small
  \caption{\textbf{Per-cell-type HD MI summary across three
  representative trained parameterisations} (one representative
  fold each: Known-ODE $+$TV, fully connected, GNN $+$TV), from a
  $T = 10\,000$ frame OU rollout ($n_\theta = 64$ bins). $\bar I_i$
  is the per-neuron mean $I(\sigma(\hat h_i);\theta)$ within the
  type; $I_{\mathrm{joint}}$ is the population MI of the whole type
  (CV-logreg lower bound). Redundancy $\rho = \sum_i I_i /
  I_{\mathrm{joint}}$: large $\rho$ means the cells of that type
  share near-identical HD tuning. Compared with Known-ODE $+$TV,
  the fully connected RNN reallocates HD onto $\Delta 7$ and ER6
  with EPG falling to $0.61$ bits, and the GNN $+$TV essentially
  silences ER6 ($0.22$ bits) while preserving the canonical
  EPG / $\Delta 7$ / PEG / EPGt backbone at a uniformly lower
  per-neuron level.}
  \label{tab:hd_mi_summary}
  \begin{tabular}{l|r|rrr|rrr|rrr}
    \toprule
    & & \multicolumn{3}{c|}{Known-ODE RNN}
      & \multicolumn{3}{c|}{fully connected RNN}
      & \multicolumn{3}{c}{GNN} \\
    type & $n$
      & $\bar I_i$ & $I_{\mathrm{joint}}$ & $\rho$
      & $\bar I_i$ & $I_{\mathrm{joint}}$ & $\rho$
      & $\bar I_i$ & $I_{\mathrm{joint}}$ & $\rho$ \\
    \midrule
    EPG    & 46 & 2.20 & 2.61 & 38.7 & 0.61 & 2.74 & 10.2 & 0.57 & 2.25 & 11.6 \\
    EPGt   &  4 & 2.33 & 1.46 &  6.4 & 0.35 & 0.86 &  1.6 & 1.16 & 1.07 &  4.3 \\
    PEN$_a$ & 20 & 1.32 & 2.71 &  9.8 & 0.41 & 2.20 &  3.7 & 0.48 & 2.18 &  4.4 \\
    PEN$_b$ & 22 & 1.00 & 2.80 &  7.8 & 0.38 & 2.27 &  3.6 & 0.62 & 2.21 &  6.1 \\
    $\Delta 7$ & 42 & 2.00 & 2.58 & 32.6 & 1.13 & 3.03 & 15.7 & 0.94 & 2.32 & 17.1 \\
    PEG    & 18 & 1.92 & 2.73 & 12.7 & 0.39 & 2.06 &  3.4 & 1.13 & 2.03 & 10.0 \\
    ER6    &  4 & 1.70 & 1.35 &  5.0 & 1.19 & 2.02 &  2.4 & 0.22 & 0.07 & 12.5 \\
    \midrule
    all (156) & --
      & \multicolumn{3}{c|}{$I_{\mathrm{pop}} = 3.18$ bits}
      & \multicolumn{3}{c|}{$I_{\mathrm{pop}} = 3.20$ bits}
      & \multicolumn{3}{c}{$I_{\mathrm{pop}} = 2.92$ bits} \\
    \bottomrule
  \end{tabular}
\end{table}

\begin{figure}[H]
  \centering
  \includegraphics[width=0.7\textwidth]{figure/fig_omega_mi_summary_nbins32.png}
  \caption{\textbf{Per-cell-type angular-velocity mutual information
  across the six parameterisations.} Same OU rollout and layout as
  \cref{fig:hd_mi_summary} with $\theta$ replaced by the instantaneous
  angular velocity $\omega$ ($n_\omega = 32$ uniform bins over the
  99\textsuperscript{th}-percentile $|\omega|$ range, empirical
  $\sigma_\omega \approx 37.46^\circ/\mathrm{s}$). Rows: Known-ODE
  $+$TV, Known-ODE no-TV, GNN $+$TV, GNN no-TV, fully connected,
  frozen $W^{\mathrm{rec}}$; left: per-neuron
  $I(\sigma(\hat h_i);\omega)$; right: joint per-type MI with the
  population joint annotated inline. The PEN$_a$ and PEN$_b$
  sub-populations are the top $\omega$-encoders in all five
  trainable parameterisations, a structural consequence of the
  four-scalar PEN velocity gate (\cref{eq:cx_pen_gate}). Population
  joint MI is $0.70$ / $0.86$ / $1.19$ / $1.21$ / $0.89$ / $0.00$
  bits for the six rows in order, giving Gaussian-channel per-step
  posteriors $\sigma_{\omega\mid\hat h} =
  \sigma_\omega\,2^{-I_{\mathrm{pop}}} \approx 23 / 21 / 16 / 16 /
  20 / 37^\circ/\mathrm{s}$ --- of the same order as the empirical
  prior in all but the GNN, so a single time step barely constrains
  $\omega$. Path integration therefore relies on accumulating this
  noisy per-step estimate into the much tighter $10\text{--}12^\circ$
  heading posterior of \cref{fig:hd_mi_summary}.}
  \label{fig:omega_mi_summary}
\end{figure}

\begin{table}[H]
  \centering
  \small
  \caption{\textbf{Per-cell-type angular-velocity MI summary},
  same rollout and definitions as \cref{tab:hd_mi_summary} but with
  $\omega$ as the target ($n_\omega = 32$ uniform bins, range
  $[-q_{99},q_{99}]$). PEN$_a$ and PEN$_b$ are the dominant
  $\omega$-encoders across the three representative trainable
  parameterisations shown (Known-ODE $+$TV, fully connected, GNN
  $+$TV; one representative training fold each). The Gaussian-channel
  per-step posterior s.d.\ $\sigma_{\omega\mid\hat h}$ is computed
  from the population joint MI as
  $\sigma_\omega\,2^{-I_{\mathrm{pop}}}$, with empirical
  $\sigma_\omega \approx 37.46^\circ/\mathrm{s}$.}
  \label{tab:omega_mi_summary}
  \begin{tabular}{l|r|rrr|rrr|rrr}
    \toprule
    & & \multicolumn{3}{c|}{Known-ODE $+$TV}
      & \multicolumn{3}{c|}{fully connected RNN}
      & \multicolumn{3}{c}{GNN $+$TV} \\
    type & $n$
      & $\bar I_i$ & $I_{\mathrm{joint}}$ & $\rho$
      & $\bar I_i$ & $I_{\mathrm{joint}}$ & $\rho$
      & $\bar I_i$ & $I_{\mathrm{joint}}$ & $\rho$ \\
    \midrule
    EPG    & 46 & 0.19 & 0.11 & 78.1 & 0.08 & 0.24 & 15.9 & 0.07 & 0.13 & 23.2 \\
    EPGt   &  4 & 0.21 & 0.01 & 84.0 & 0.06 & 0.01 & 30.5 & 0.25 & 0.09 & 10.6 \\
    PEN$_a$ & 20 & 0.26 & 0.54 &  9.5 & 0.10 & 0.49 &  3.9 & 0.74 & 0.96 & 15.5 \\
    PEN$_b$ & 22 & 0.26 & 0.59 &  9.7 & 0.13 & 0.69 &  4.2 & 0.55 & 0.90 & 13.4 \\
    $\Delta 7$ & 42 & 0.20 & 0.09 & 89.5 & 0.15 & 0.37 & 16.7 & 0.11 & 0.21 & 21.6 \\
    PEG    & 18 & 0.19 & 0.12 & 28.0 & 0.06 & 0.11 &  9.7 & 0.15 & 0.15 & 17.8 \\
    ER6    &  4 & 0.14 & 0.04 & 15.7 & 0.12 & 0.02 & 22.5 & 0.18 & 0.14 &  5.2 \\
    \midrule
    all (156) & --
      & \multicolumn{3}{c|}{$I_{\mathrm{pop}} = 0.70$;\,$\sigma_{\omega\mid\hat h} = 23.0^\circ$/s}
      & \multicolumn{3}{c|}{$I_{\mathrm{pop}} = 0.89$;\,$20.2^\circ$/s}
      & \multicolumn{3}{c}{$I_{\mathrm{pop}} = 1.19$;\,$16.4^\circ$/s} \\
    \bottomrule
  \end{tabular}
\end{table}

\paragraph{Four functional classes of CX neurons.}
\label{para:four_classes}
Combining the per-neuron HD MI of \cref{fig:hd_mi_summary} with
the readout magnitude and the angular-velocity modulation, we apply
the R/L/D/Z partition of \cref{method:four_classes} to each
trained network. Per-cell-type counts are summarised in
\cref{tab:cx_four_classes} and the underlying scatters /
stacked-bar in \cref{fig:cx_four_classes}.

Several factual patterns emerge under the EPG-only readout.
\textbf{(i)}~$\cls{R}$ is exclusively populated by EPG neurons
(all R-class cells are EPG under both Known-ODE no-TV and GNN
no-TV; $0$ non-EPG cells in $\cls{R}$ under either model). This
is a structural consequence of $\hat W^{\mathrm{out}}_{:,i}
\equiv 0$ for $i \ge 46$: only EPG cells can simultaneously
satisfy high $\mi$ and high $\omag$, since the decoder has no
non-EPG columns to push $\omag$ onto.
\textbf{(ii)}~$\cls{D}$ is dominated by the PEN sub-populations
($\sim 70\%$ of $\cls{D}$ under either trainable architecture).
This is consistent with PEN's anatomical role as the
angular-velocity input arm of the ring-attractor: PEN$_a$ and
PEN$_b$ are the only cell types whose stimulus row carries
$\omega$ mass (\cref{method:RNN}).
\textbf{(iii)}~$\cls{L}$ is the largest class ($\sim\!60$ cells
in both nominal recipes) and gathers the HD-encoding non-EPG
residue plus the EPG cells that do not sit in the top of the
$\omag$ distribution. The $\Delta 7$ population sits predominantly
in $\cls{L}$ under both models, consistent with $\Delta 7$
acquiring HD information through the recurrent ring while being
invisible to the EPG-only readout.
\textbf{(iv)}~$\cls{Z}$ is the smallest class and is mostly
populated by EPG cells whose decoder-column magnitude and HD-MI
both sit below the medians, plus a few $\Delta 7$ / PEG / ER6
cells. Across the four trainable parameterisations the gross
structure of the partition (EPG carries $\cls{R}\cup\cls{L}$;
PEN carries $\cls{D}$; $\Delta 7$ / PEG carry $\cls{L}$; ER6 /
EPGt sit in small fractions of $\cls{L}$ or $\cls{Z}$) is
preserved, so the same circuit algorithm is recovered up to a
quantitative shuffling of borderline cells.

\begin{table}[H]
  \centering
  \small
  \caption{\textbf{Per-cell-type breakdown across the four
  functional classes} for the two nominal no-TV recipes (Known-ODE
  no-TV and GNN no-TV), one representative training fold each.
  Class definitions: $\cls{R}$ high-MI \& high $\omag$, $\cls{L}$
  high-MI \& low $\omag$, $\cls{D}$ low-MI \& high $\omod$,
  $\cls{Z}$ low-MI \& low $\omod$; thresholds set at the sample
  median (taken over the non-zero $\omag$ entries for the $\omag$
  threshold). Class names follow the convention of the zebrafish
  HD-circuit partition \citep{petrucco2023}.}
  \label{tab:cx_four_classes}
  \begin{tabular}{l rrrrr rrrrr}
\toprule
type & \multicolumn{5}{c}{Known-ODE no-TV} & \multicolumn{5}{c}{GNN no-TV} \\
\midrule
 & R & L & D & Z & total & R & L & D & Z & total \\
\midrule
EPG & 11 & 10 & 16 & 9 & 46 & 13 & 15 & 6 & 12 & 46 \\
EPGt & 0 & 3 & 0 & 1 & 4 & 0 & 0 & 3 & 1 & 4 \\
PEN$_a$ & 0 & 0 & 20 & 0 & 20 & 0 & 3 & 14 & 3 & 20 \\
PEN$_b$ & 0 & 0 & 22 & 0 & 22 & 0 & 0 & 21 & 1 & 22 \\
PEG & 0 & 15 & 0 & 3 & 18 & 0 & 9 & 4 & 5 & 18 \\
$\Delta 7$ & 0 & 39 & 3 & 0 & 42 & 0 & 36 & 4 & 2 & 42 \\
ER6 & 0 & 0 & 0 & 4 & 4 & 0 & 2 & 1 & 1 & 4 \\
\midrule
total & 11 & 67 & 61 & 17 & 156 & 13 & 65 & 53 & 25 & 156 \\
\bottomrule
\end{tabular}

\end{table}

\begin{figure}[H]
  \centering
  \includegraphics[width=\textwidth]{figure/fig_cx_four_classes.png}
  \caption{\textbf{Four functional classes of CX neurons under
  the EPG-only readout.} Columns: Known-ODE no-TV, Known-ODE $+$TV,
  GNN no-TV, GNN $+$TV (one representative training fold each).
  \textbf{Row 1:}~heading mutual information $\mi$ versus readout
  magnitude $\omag_i$, with median thresholds (dashed). The
  $\cls{R}$ vs.\ $\cls{L}$ split is set by $\omag^\star$.
  \textbf{Row 2:}~MI versus angular-velocity modulation $\omod_i$;
  the $\cls{D}$ vs.\ $\cls{Z}$ split sits at $\omod^\star$.
  \textbf{Row 3:}~per-cell-type stacked counts for the seven
  CX cell types ($\cls{R}$ red, $\cls{L}$ orange, $\cls{D}$ blue,
  $\cls{Z}$ grey). EPG cells exclusively populate $\cls{R}$ as
  a structural consequence of the EPG-only decoder
  ($\hat W^{\mathrm{out}}_{:,i} \equiv 0$ for $i \ge n_{\mathrm{EPG}}
  = 46$), not as an emergent property of the trained dynamics:
  only EPG columns of $\hat W^{\mathrm{out}}$ are non-zero, so only
  EPG cells can satisfy high $\omag$. PEN$_a$ / PEN$_b$ supply
  $\sim 70\%$ of $\cls{D}$; $\Delta 7$, PEG and EPGt sit
  predominantly in $\cls{L}$. The gross structure is preserved
  across the four trainable parameterisations
  (cf.\ \cref{tab:cx_four_classes}).}
  \label{fig:cx_four_classes}
\end{figure}


\clearpage

\section*{Appendix A: Degeneracy of the inverse problem}
\label{app:degeneracy_cx}

Above, we demonstrated that the trained recurrent neural network with imposed connectome topology and Dale signs reproduces the path-integration task at near-perfect accuracy across 10 independently trained CV folds (\cref{fig:evolution_pi}, \cref{fig:w_rec_comparison}). Here, we ask how degenerate the inverse problem is: how many distinct recurrent weight matrices $\hat{W}^{\mathrm{rec}}$ reproduce the same voltage trajectories, and which directions in $\hat{W}^{\mathrm{rec}}$-space are actually invariant for the trained solutions? All numerical results in this appendix are reported as mean $\pm$ std across the same 10 CV folds used in \cref{fig:w_rec_comparison}.

We fix the values of the bias $\hat{b}$, the time constant $\tau$, and the input weights $\hat{W}^{\mathrm{in}}$ at their trained values, and rearrange \cref{eq:cx_dynamics} to isolate the incoming synaptic weights:
\begin{equation}
\sum_{j\in\mathcal{N}_i} \hat{W}^{\mathrm{rec}}_{ij}\,\sigma\!\big(\hat{h}_j(t)\big)
\;=\; \tau\,\frac{d\hat{h}_i(t)}{dt} + \hat{h}_i(t) - \hat{b}_i - I_i(t),
\end{equation}
where $I_i(t) = \sum_l \hat{W}^{\mathrm{in}}_{il}\,u_l(t)$ is the per-neuron projected input from \cref{eq:cx_pen_gate}. Defining the presynaptic gated activity $h_j(t) = \sigma(\hat{h}_j(t))$ and the residual $b_i(t) = \tau\,d\hat{h}_i(t)/dt + \hat{h}_i(t) - \hat{b}_i - I_i(t)$, and stacking the $T$ time-point constraints into matrix form, yields an independent linear system for each postsynaptic neuron $i$:
\begin{equation}
\mathbf{H}_i\,\mathbf{w}_i = \mathbf{b}_i,
\label{eq:cx_linear_system}
\end{equation}
where $\mathbf{H}_i \in \mathbb{R}^{T \times d_i}$ is the presynaptic activity matrix, $d_i$ is the in-degree of neuron $i$, $\mathbf{w}_i \in \mathbb{R}^{d_i}$ the incoming recurrent weight vector, and $\mathbf{b}_i \in \mathbb{R}^T$ the stacked residuals. This holds for each postsynaptic neuron $i = 1, \ldots, N$, where $N = 156$ and $d_i$ ranges from $1$ to $98$. The null space $\ker(\mathbf{H}_i) = \{\boldsymbol{\delta} \in \mathbb{R}^{d_i} : \mathbf{H}_i\boldsymbol{\delta} = \mathbf{0}\}$ characterises the degeneracy: any $\mathbf{w}_i + \boldsymbol{\delta}$ with $\boldsymbol{\delta} \in \ker(\mathbf{H}_i)$ satisfies \cref{eq:cx_linear_system} equally and produces identical dynamics. In the following, we characterise the dynamical equivalence class of $\hat{W}^{\mathrm{rec}}$, the set of all recurrent weight matrices that reproduce the trained model's voltage trajectories.

\paragraph{Two SVD-based bounds on the null space (inconclusive).}
We solve \cref{eq:cx_linear_system} per neuron via column-equilibrated normal equations and the eigendecomposition of $\mathbf{H}_i^{\top}\mathbf{H}_i$, and quantify $\dim\ker(\mathbf{H})$ in two complementary ways. \emph{Strict per-neuron SVD:} the eigendecomposition flags directions as \emph{null} (relative eigenvalue $\le T \cdot 10^{-22}$) or \emph{sloppy} ($\le T \cdot 10^{-12}$); summed over $i$ this counts only $6.3\%$ of edges ($0$ null + $650$ sloppy). \emph{Global-SVD coarse bound:} by rank--nullity, $d_i = \mathrm{rank}(\mathbf{H}_i) + \dim\ker(\mathbf{H}_i)$. Each per-neuron activity matrix $\mathbf{H}_i$ is a column-submatrix of the global activity matrix $\mathbf{H} \in \mathbb{R}^{T \times N}$ stacking sigmoid gated activities $\sigma(\hat h_j(t))$ of all $156$ neurons over $T = 8{,}000$ subsampled frames; therefore $\mathrm{rank}(\mathbf{H}_i) \le r$, where $r$ is the effective rank of $\mathbf{H}$ at a chosen variance threshold, and
\begin{equation}
\dim\ker(\mathbf{H}) \;\approx\; \sum_{i=1}^{N} \max\!\big(0,\; d_i - r\big).
\label{eq:cx_null_globalsvd}
\end{equation}
At $99\%$ variance, $r = 6$ --- only six modes carry $99\%$ of the population variance, reflecting the low-dimensional ring-attractor manifold --- and the bound spans $9{,}327$ edges ($90.9\%$). Tightening the threshold to $99.9\%$ raises $r$ to $15$ and shrinks the bound to $7{,}923$ edges ($77.2\%$); the null space remains large at every threshold. The two estimates disagree by more than an order of magnitude and neither is conclusive on its own. To pin down which directions are actually invariant for the trained model, we apply controlled perturbations to specific groups of edges and roll out the network. We test two partition schemes in the next two paragraphs.

\paragraph{Per cell type.}
For each postsynaptic neuron $i$ and presynaptic cell type $\alpha$, the $k_{i,\alpha}$ same-type incoming edges form a candidate $(i,\alpha)$ group; pairs with $k_{i,\alpha} \ge 2$ account for $9{,}400$ edges ($91.6\%$ of the $10{,}263$-edge connectome support; structural, identical across folds). The sum-zero perturbation is the natural test for this partition: when the $k_{i,\alpha}$ same-type presynaptic neurons carry correlated subthreshold activity, the recurrent drive $\sum_m W_{i,j_m}\,\sigma(\hat h_{j_m}(t)) \approx \sigma(\hat h(t))\,\sum_m W_{i,j_m}$ depends only on the group sum, so any perturbation on the same-type weights $W_{i,j_m} \to W_{i,j_m} + \delta_{j_m}$ with $\sum_m \delta_{j_m} = 0$ leaves the dynamics invariant under the lockstep assumption \citep{allier_connectome_2026}. We test this candidate null direction empirically by perturbing each fold's groups by a sum-zero vector at amplitude $\lambda \cdot \mathrm{mean}(|W_{\mathrm{group}}|)$ and rolling out for $1{,}000$ frames against that fold's own GT voltage; reported values average over the 10 folds (\cref{tab:variants_cx}). At $\lambda = 1$, all seven cell types preserve the dynamics ($R^{2}_{\mathbf{W}} \in [0.97, 1.00]$ across folds, Pearson $r$ in $[0.96, 1.00]$): the predicted directions are real \emph{locally} around the trained $\hat{W}^{\mathrm{rec}}$, and the result is reproducible across folds (per-cell-type $r$ SDs $\le 0.09$). At $\lambda = 2$ all $r$ remain above $0.88$ (PEG and $\Delta 7$ the loosest, the latter with the largest inter-fold spread, $r = 0.88 \pm 0.14$); at $\lambda = 4$, $R^{2}_{\mathbf{W}}$ drops to $[0.54, 0.97]$ and $\Delta 7$/PEG rollouts deteriorate most strongly ($r = 0.83 \pm 0.08$ / $0.86 \pm 0.12$). Pushing all the way to the extreme by collapsing each $(i,\alpha)$ group of $k_{i,\alpha}$ edges onto a single representative edge $e^{\star} = \arg\max_e |W_e|$ --- setting $W_{e^{\star}} = \sum_{e \in \mathrm{group}} W_e$ (sum-preserving) or $W_{e^{\star}}$ chosen by a one-scalar least-squares fit that minimises the per-group dynamics residual (calibrated), and zeroing the other $k_{i,\alpha} - 1$ edges --- removes $91.6\%$ of the edges. This is the analogue of the operation that preserves dynamics at $r \approx 0.95$ on the visual-system biomodel of \citep{allier_connectome_2026}, but it breaks the CX heading readout for both variants: sum-preserving rollout $r = 0.27 \pm 0.24$, calibrated rollout $r = 0.13 \pm 0.18$ (both averaged over 10 folds). The per-cell-type null space is therefore approximately null around the trained weights but cannot be saturated by a single representative per cell-type pair.

\paragraph{Per unit.}
The per-$(i,\alpha)$ partition lumps same-type neurons regardless of which PB glomerulus they innervate, but in the CX two EPGs in different glomeruli encode different headings and are not clones. The natural refinement groups same-type neurons by their hemibrain instance, the per-cell annotation that identifies a neuron's anatomical computational unit: for EPG/EPGt/PEG/PEN it is the PB glomerulus (e.g., \emph{EPG(PB08)\_L3} marks an EPG neuron in glomerulus L3); for $\Delta 7$ it is the set of glomeruli the neuron spans (e.g., \emph{Delta7(PB15)\_L1L9R8\_R}); for ER6 it is the hemisphere. Same-type neurons sharing an instance are clones as discussed in \citet{hulse2026}. The $(i,\alpha,\mathrm{instance})$ partition retains $5{,}150$ candidate edges ($50.2\%$ of $10{,}263$; structural, identical across folds). Applying the same collapse defined above to each $(i,\alpha,\mathrm{instance})$ group --- zeroing $50.2\%$ of the edges instead of $91.6\%$ --- yields, across 10 folds, sum-preserving rollout $r = 0.55 \pm 0.20$ and calibrated rollout $r = 0.25 \pm 0.28$, against $r = 0.27 \pm 0.24$ / $0.13 \pm 0.18$ for the analogous operations under the per-$(i,\alpha)$ partition. The per-unit sum-preserving collapse therefore continues to track heading more robustly than the per-cell-type one, confirming that the finer (instance-level) partition more closely respects the trained model's null space. The single-fold $r = 0.98$ calibrated value reported in the predecessor (no-TV-prior) analysis no longer holds once $\hat W^{\mathrm{rec}}$ is biased away from the per-instance null directions by the TV regulariser used in the present training recipe (\cref{para:agent_gnn}): across folds, the calibrated collapse falls below the sum-preserving one and has large inter-fold spread ($0.25 \pm 0.28$, 95\% CI spans the full $[0, 0.8]$ range). Mechanistically, this says that the trained $\hat W^{\mathrm{rec}}$ is not symmetric within an $(\alpha, \mathrm{instance})$ clone group: each fold finds a structurally similar but non-interchangeable per-clone allocation, and a single-scalar OLS fit cannot reproduce that asymmetry as accurately as the structural sum. In a retinotopic visual system \citep{allier_connectome_2026}, the $(i,\alpha)$ partition already saturates because columns of the same type encode the same scene; the CX requires the finer $(i,\alpha,\mathrm{instance})$ partition because heading is encoded per-glomerulus, not per-type --- but on the trained-weight side, the calibrated saturation that holds for the visual system is broken here by the TV-shaped solutions.

\begin{figure}[H]
\centering
\includegraphics[width=\textwidth]{figure/fig_sparsify_cx.png}
\caption{\textbf{Sparsifying $\hat W^{\mathrm{rec}}$ along the candidate null directions.} Top row (a-d): recurrent matrices, sorted by cell type (block boundaries marked) and z-scored on the same scale (computed from the GT non-zero entries, clipped to $\pm 3$). Bottom row (e-h): heading on one OU test trial decoded from each rollout; true HD in light green, decoded HD in black, both wrapped to $(-\pi, \pi]$. \textbf{(a, e)} Learned $\hat W^{\mathrm{rec}}$ from the trained CxTaskRNN: GT and decoded HD coincide. \textbf{(b, f)} Per-$(i,\alpha)$ sum-preserving collapse, $91.6\%$ of edges zeroed: the bump-coding blocks (EPG, PEN$_a$, PEN$_b$, $\Delta 7$) lose most of their structure and the decoded HD drifts away from GT ($r = 0.04$). \textbf{(c, g)} Per-$(i,\alpha,\mathrm{instance})$ sum-preserving collapse, $50.2\%$ of edges zeroed: the block structure is preserved and the decoded HD partially tracks GT ($r = 0.53$). \textbf{(d, h)} Same per-$(i,\alpha,\mathrm{instance})$ partition but with a \emph{calibrated} one-scalar-per-group fit (each group of $k$ same-instance clones is replaced by a single representative edge whose weight is set by least-squares to minimise the per-group dynamics residual): the same $50.2\%$ of edges are zeroed yet the bump tracks GT at $r = 0.98$, demonstrating that the per-$(i,\alpha,\mathrm{instance})$ partition is near-fully saturable when intra-unit clones are summarised by a calibrated scalar instead of a naïve sum.}
\label{fig:sparsify_cx}
\end{figure}

\begin{table}[H]
\centering
\footnotesize
\begin{tabular}{lrrrrrr}
\toprule
Cell Type & \multicolumn{2}{c}{$\lambda = 1$} & \multicolumn{2}{c}{$\lambda = 2$} & \multicolumn{2}{c}{$\lambda = 4$} \\
 & $R^2_{\mathbf{W}}$ & $r$ & $R^2_{\mathbf{W}}$ & $r$ & $R^2_{\mathbf{W}}$ & $r$ \\
\midrule
EPG & $0.998\pm0.000$ & $0.999\pm0.002$ & $0.992\pm0.002$ & $0.997\pm0.003$ & $0.970\pm0.007$ & $0.975\pm0.034$ \\
EPGt & $0.998\pm0.001$ & $0.990\pm0.029$ & $0.992\pm0.003$ & $0.953\pm0.095$ & $0.969\pm0.012$ & $0.949\pm0.115$ \\
PEN\_a(PEN1) & $0.994\pm0.001$ & $0.986\pm0.028$ & $0.974\pm0.005$ & $0.991\pm0.013$ & $0.898\pm0.022$ & $0.936\pm0.066$ \\
PEN\_b(PEN2) & $0.988\pm0.003$ & $0.964\pm0.088$ & $0.951\pm0.012$ & $0.981\pm0.016$ & $0.803\pm0.050$ & $0.940\pm0.088$ \\
Delta7 & $0.971\pm0.004$ & $0.960\pm0.049$ & $0.886\pm0.015$ & $0.876\pm0.142$ & $0.543\pm0.059$ & $0.833\pm0.078$ \\
PEG & $0.989\pm0.003$ & $0.992\pm0.006$ & $0.957\pm0.012$ & $0.979\pm0.017$ & $0.826\pm0.050$ & $0.859\pm0.117$ \\
ER6 & $0.991\pm0.002$ & $0.999\pm0.002$ & $0.963\pm0.006$ & $0.990\pm0.020$ & $0.851\pm0.025$ & $0.985\pm0.028$ \\
\bottomrule
\end{tabular}

\caption{\textbf{Single-type sum-zero variants across perturbation amplitudes $\lambda$, mean $\pm$ std over 10 CV folds.} For each cell type with $k_{i,\alpha} \ge 2$ groups, all degenerate weights are perturbed by a sum-zero vector at amplitude $\lambda\cdot\mathrm{mean}(|W_{\mathrm{group}}|)$ and rolled out for $1{,}000$ frames against that fold's own GT voltage trajectory. Columns: connectivity $R^{2}_{\mathbf{W}}$ (between the perturbed and trained recurrent matrix on the connectome support) and rollout Pearson $r$ at $\lambda \in \{1, 2, 4\}$, averaged across the 10 CV folds with the inter-fold standard deviation as $\pm$ band. At $\lambda = 1$ every type is preserved with $r \ge 0.96 \pm 0.09$; both $R^{2}_{\mathbf{W}}$ and $r$ decrease with $\lambda$, confirming that the candidate null directions are only approximately null at large perturbations, and the inter-fold spread grows monotonically with $\lambda$.}
\label{tab:variants_cx}
\end{table}

\clearpage

\section*{Appendix B: Training hyperparameters}
\label{app:hyperparams}

\Cref{tab:cv_hyperparams} lists the training hyperparameters used
for the cv0 fold of each of the six conditions; the remaining nine
folds in every condition share the same configuration and differ
only in their random seed (simulation, training, and task).
$\eta_{\mathrm{rec}}$ is the learning rate applied to the recurrent
weights ($\hat S$ for the Known-ODE and fully connected RNNs, the
GNN MLPs and embeddings for the GNN, and zero for the frozen
$W^{\mathrm{rec}}$ control). $\eta_{\mathrm{enc/dec}}$ is the
learning rate applied to the encoder $\hat W^{\mathrm{in}}$ /
decoder $\hat W^{\mathrm{out}}$ blocks (including the four PEN
angular-velocity scalars), kept distinct because these blocks are
small and converge much faster than $\hat S$.
$\lambda_{\mathrm{tv}}$ is the EPG-ring TV coefficient of
\cref{eq:cx_tv}, set to zero in the no-TV conditions.
$\lambda_{\mathrm{nf}}$ is the connectome-block norm-floor
coefficient (active only on the fully connected RNN to prevent the
column-wise Dale split from collapsing). $\epsilon_{\mathrm{tail}}$
is the tail-loss coefficient on the decoded heading squared error
that anchors the readout at training time. $\mu_f / \mu_g$ are the
GNN monotonicity priors on $f_\theta$ and $g_\phi$ (leak-monotone
in $\hat h$, drive-monotone in the aggregated input $m$). The
remaining hyperparameters not listed in the table are shared
across conditions: the integration step $\Delta t = 0.01$~s, the
training rollout length $T = 1000$ frames, the AdamW optimiser, a
constant learning-rate schedule, and the cos/sin training
target of \cref{eq:cx_readout}. The cosine-distance and
$\hat W^{\mathrm{rec}}$ $\ell_1$ coefficients
($\lambda_{\mathrm{cos}}$, $\lambda_{W,1}$) were zero throughout
this study.

\begin{table}[H]
\centering
\small
% Training hyperparameters per condition (10 CV folds each).
% Columns: condition | lr (rec) | lr (in/out, encoder/decoder) | lambda_tv | lambda_nf | epsilon (tail-loss) | mu_f, mu_g (GNN drift / message monotonicity) | n_epochs | batch
\begin{tabular}{lrrrrrrrr}
\toprule
condition & $\eta_{\mathrm{rec}}$ & $\eta_{\mathrm{enc/dec}}$
  & $\lambda_{\mathrm{tv}}$ & $\lambda_{\mathrm{nf}}$ & $\epsilon_{\mathrm{tail}}$
  & $\mu_{f}/\mu_{g}$ & $n_{\mathrm{epochs}}$ & $|\mathcal B|$ \\
\midrule
Known-ODE $+$TV         & $2\!\times\!10^{-3}$ & $2.5\!\times\!10^{-3}$ & $10^{-2}$ & $0$    & $3.5\!\times\!10^{-2}$ & --              & 10 & 64 \\
Known-ODE no-TV         & $2\!\times\!10^{-3}$ & $2.5\!\times\!10^{-3}$ & $0$       & $0$    & $3.5\!\times\!10^{-2}$ & --              & 10 & 64 \\
GNN $+$TV               & $2\!\times\!10^{-3}$ & $5\!\times\!10^{-4}$   & $5\!\times\!10^{-3}$ & $0$ & $5\!\times\!10^{-2}$ & $15/15$ & 10 &  8 \\
GNN no-TV               & $2\!\times\!10^{-3}$ & $5\!\times\!10^{-4}$   & $0$       & $0$    & $5\!\times\!10^{-2}$ & $15/15$ & 10 &  8 \\
fully connected         & $2\!\times\!10^{-3}$ & $2.5\!\times\!10^{-3}$ & $0$       & $0.5$  & $5\!\times\!10^{-2}$ & --              &  5 & 64 \\
frozen $W^{\mathrm{rec}}$ & $0$                & $6\!\times\!10^{-3}$   & $0$       & $0$    & $5\!\times\!10^{-2}$ & --              &  1 & 64 \\
\bottomrule
\end{tabular}

\caption{\textbf{Training hyperparameters for the cv0 fold of each
of the six conditions} (the other nine folds share the same
configuration). $\eta_{\mathrm{rec}}$, $\eta_{\mathrm{enc/dec}}$:
learning rates on the recurrent / encoder--decoder parameter
blocks. $\lambda_{\mathrm{tv}}$, $\lambda_{\mathrm{nf}}$:
EPG-ring TV (\cref{eq:cx_tv}) and per-block norm-floor coefficients.
$\epsilon_{\mathrm{tail}}$: tail-loss coefficient on the decoded
heading. $\mu_f/\mu_g$: GNN drift / message monotonicity priors
of \cref{eq:cx_gnn_dynamics}. $n_{\mathrm{epochs}}$:
training epochs. $|\mathcal B|$: batch size.
``--'' marks coefficients that do not exist in the corresponding
architecture (e.g.\ $\mu_f, \mu_g$ for the Known-ODE / FC / frozen
RNNs).}
\label{tab:cv_hyperparams}
\end{table}

\clearpage

\begin{thebibliography}{9}

\bibitem[Allier et~al.(2026a)]{allier_connectome_2026}
C.~Allier et~al.
\newblock Graph neural networks recover interpretable circuit models
  from neural activity.
\newblock In preparation (NeurIPS 2026 submission), 2026.

\bibitem[Allier and Saalfeld(2026b)]{allier_gnn_transformer_2026}
C.~Allier and S.~Saalfeld.
\newblock A four-MLP GNN framework with explicit temporal filtering, and
  the cost-structured route to emergent circuit architecture.
\newblock Internal technical report, 2026.

\bibitem[Gilmer et~al.(2017)]{gilmer_neural_2017}
J.~Gilmer, S.~S.~Schoenholz, P.~F.~Riley, O.~Vinyals, and G.~E.~Dahl.
\newblock Neural message passing for quantum chemistry.
\newblock In \emph{ICML}, 2017.

\bibitem[Hulse et~al.(2026)]{hulse2026}
B.~K. Hulse, P.~B. Aneesh, S.~Romani, V.~Jayaraman, and A.~M. Hermundstad.
\newblock How hidden symmetries in network connectivity support robust
  attractor dynamics in the fly's neural compass.
\newblock \emph{bioRxiv}, 2026.
  doi:\,\href{https://doi.org/10.64898/2026.05.18.725766}{10.64898/2026.05.18.725766}.
\newblock URL \url{https://www.biorxiv.org/content/10.64898/2026.05.18.725766v1}.

\bibitem[Green et~al.(2017)]{green2017}
J.~Green, A.~Adachi, K.~K. Shah, J.~D. Hirokawa, P.~S. Magani, and
  G.~Maimon.
\newblock A neural circuit architecture for angular integration in
  \emph{Drosophila}.
\newblock \emph{Nature}, 546(7656):101--106, 2017.

\bibitem[Hulse et~al.(2021)]{hulse2021}
B.~K. Hulse et~al.
\newblock A connectome of the \emph{Drosophila} central complex reveals
  network motifs suitable for flexible navigation and context-dependent
  action selection.
\newblock \emph{eLife}, 10:e66039, 2021.

\bibitem[Kim et~al.(2017)]{kim2017}
S.~S. Kim, H.~Rouault, S.~Druckmann, and V.~Jayaraman.
\newblock Ring attractor dynamics in the \emph{Drosophila} central brain.
\newblock \emph{Science}, 356(6340):849--853, 2017.

\bibitem[Seelig and Jayaraman(2015)]{seelig2015}
J.~D. Seelig and V.~Jayaraman.
\newblock Neural dynamics for landmark orientation and angular path
  integration.
\newblock \emph{Nature}, 521(7551):186--191, 2015.

\bibitem[Vafidis et~al.(2022)]{vafidis2022}
P.~Vafidis, D.~Owald, T.~D'Albis, and R.~Kempter.
\newblock Learning accurate path integration in ring attractor models
  of the head direction system.
\newblock \emph{eLife}, 11:e69841, 2022.

\bibitem[Petrucco et~al.(2023)]{petrucco2023}
L.~Petrucco, H.~Lavian, Y.~K. Wu, F.~Svara, V.~\v{S}tih, R.~Portugues.
\newblock Neural dynamics and architecture of the heading direction
  circuit in zebrafish.
\newblock \emph{Nature Neuroscience}, 26:765--773, 2023.

\end{thebibliography}

\end{document}